%~Mouliné par MaN_auto v.0.27.3 2023-05-26 11:24:32
\documentclass[AHL,Unicode,longabstracts,published]{cedram}

\usepackage{bbm}
\usepackage{mathrsfs}
\usepackage{CJK}








%\theoremstyle{plain}


\newcommand{\ddr}{\mathrm{d}}
\newcommand{\rmr}{\mathrm{r}}
\newcommand{\rmm}{\mathrm{m}}
\newcommand{\rmk}{\mathrm{k}}
\newcommand{\rms}{\mathrm{s}}
\newcommand{\PPP}{\mathrm{PPP}}
\newcommand{\MP}{\mathrm{MP}}
\newcommand{\Frag}{\mathrm{Frag}}
\newcommand{\Ber}{\mathrm{Ber}}
\newcommand{\bbmun}{\mathbbm{1}}
\newcommand{\bbN}{\mathbb{N}}
\newcommand{\bbR}{\mathbb{R}}
\newcommand{\bbE}{\mathbb{E}}
\newcommand{\bbP}{\mathbb{P}}
\newcommand{\bbH}{\mathbb{H}}
\newcommand{\clM}{\mathcal{M}}
\newcommand{\clH}{\mathcal{H}}
\newcommand{\clA}{\mathcal{A}}
\newcommand{\clL}{\mathcal{L}}
\newcommand{\clP}{\mathcal{P}}

\newcounter{stepcount}
%\renewcommand{\thestepcount}{\alph{stepcount}}
\newenvironment{step}{\refstepcounter{stepcount}\begin{proof}[\textbf{Step~\thestepcount}]}{\end{proof}}

 
%%%%%%%%%%%%%%%%%%%%%%%%%%%%%%%%%%%%%%%%%%%%%%%%%%%%%%%%%%%%%%%%
\graphicspath{{./figures/}}

\newcommand*{\mk}{\mkern -1mu}
\newcommand*{\Mk}{\mkern -2mu}
\newcommand*{\mK}{\mkern 1mu}
\newcommand*{\MK}{\mkern 2mu}

%\hypersetup{urlcolor=purple, linkcolor=blue, citecolor=red}

\newcommand*{\romanenumi}{\renewcommand*{\theenumi}{\roman{enumi}}}
\newcommand*{\Romanenumi}{\renewcommand*{\theenumi}{\Roman{enumi}}}
\newcommand*{\alphenumi}{\renewcommand*{\theenumi}{\alph{enumi}}}
\newcommand*{\Alphenumi}{\renewcommand*{\theenumi}{\Alph{enumi}}}
\let\oldtilde\tilde
\renewcommand*{\tilde}[1]{\mathchoice{\widetilde{#1}}{\widetilde{#1}}{\oldtilde{#1}}{\oldtilde{#1}}}
\let\oldexists\exists
\renewcommand*{\exists}{\mathrel{\oldexists}}
\let\oldforall\forall
\renewcommand*{\forall}{\mathrel{\oldforall}}
%%%%%%%%%%%%%%%%%%%%%%%%%%%%%%%%%%%%%%%%%%%%%%%%%%%%%%%%%%%%%%%

\title[$\Lambda$-Wright--Fisher processes with selection \& simple EFC processes]{On the boundary classification of $\Lambda$-Wright--Fisher processes with frequency-dependent selection}
\alttitle{Sur la classification des points frontières des processus de $\Lambda$-Wright--Fisher avec sélection dépendante de la fréquence}

\subjclass{60J80, 60J70, 60J90, 92D25}
\keywords{$\Lambda$-Wright--Fisher process, selection, $\Lambda$-coalescent, fragmentation-coalescence, duality, explosion, coming down from infinity, entrance boundary, regular boundary, continuous-time Markov chains}

\author[\initial{C.} \lastname{Foucart}]{\firstname{Clément} \lastname{Foucart}}
\address{Université Sorbonne\\
Paris Nord and Paris 8,\\
Laboratoire Analyse,\\
Géométrie \& Applications,\\
UMR 7539.\\
Institut Galilée,\\
99 avenue J.B. Clément,\\
93430 Villetaneuse, France}
\email{foucart@math.univ-paris13.fr}
\thanks{C.F's research is partially supported by LABEX MME-DII (ANR11-LBX-0023-01). X.Z.'s research is supported by Natural Sciences and Engineering Research Council of Canada (RGPIN-2016-06704) and by National Natural Science Foundation of China (No.\ 11731012).}


\author[\initial{X.} \lastname{Zhou}]{\firstname{Xiaowen} \lastname{Zhou}}
\address{Department of Mathematics\\
and Statistics,\\
Concordia University,\\
1455 De Maisonneuve Blvd. W.,\\
Montreal, Canada}
\email{xiaowen.zhou@concordia.ca}





\begin{abstract}
We construct extensions of the pure-jump $\Lambda$-Wright--Fisher processes with frequency-dependent selection~($\Lambda$-WF with selection) with different behaviors at their boundary $1$. Those processes satisfy some duality relationships with the block counting process of simple exchangeable fragmentation-coagulation processes (EFC processes). One-to-one correspondences are established between the nature of the boundaries $1$ and $\infty$ of the processes involved. They provide new information on these two classes of processes. Sufficient conditions are provided for boundary $1$ to be an exit boundary or an entrance boundary. When the coalescence measure $\Lambda$ and the selection mechanism verify some regular variation properties, conditions are found in order that the extended $\Lambda$-WF process with selection makes excursions out from the boundary $1$ before getting absorbed at $0$. In this case, $1$ is a transient regular reflecting boundary. This corresponds to a new phenomenon for the deleterious allele, which can be carried by the whole population for a set of times of zero Lebesgue measure, before vanishing in finite time almost surely.
\end{abstract}

\begin{altabstract}
Nous construisons des extensions des processus de $\Lambda$-Wright--Fisher de saut pur avec sélection dépendante de la fréquence ($\Lambda$-WF avec sélection) présentant différents comportement en leur point frontière $1$. Ces processus satisfont des relations de dualité avec le processus du nombre de blocs des processus de fragmentation-coagulation échangeables simples. Des correspondances biunivoques entre les natures des frontières $1$ et $\infty$ des processus en question sont établies. Elles fournissent de nouvelles informations sur ces deux classes de processus. Des conditions suffisantes sont données pour que la frontière $1$ soit un point de sortie ou un point d'entrée. Lorsque la mesure $\Lambda$ et la fonction de sélection~vérifient des propriétés de variations régulières, des conditions sont trouvées de sorte que le processus de $\Lambda$-WF avec sélection étendu fasse des excursions en dehors de la frontière $1$ avant d'être absorbée en $0$. Dans ce cas, $1$ est un point régulier réfléchissant et transient. Cela correspond à un nouveau phénomène pour l'allèle délétère, qui peut être porté par toute la population pendant un ensemble de temps de mesure de Lebesgue nulle, avant de disparaître en temps fini presque sûrement.
\end{altabstract}


\datereceived{2021-01-07}
\daterevised{2021-05-12}
\dateaccepted{2021-09-14}

\editor{L. Chaumont}

%%%%%%%%%%%%%%%%%%%%%%%%%%%%%%%%%%%%%%%%%%%%%%%%%%%%%%%%%%%%%%%%%	
\dateposted{2023-07-18}
\begin{document}

\maketitle
\section{Introduction}
Wright--Fisher processes are fundamental mathematical models in population genetics. They are Markov processes, taking their values in the interval $[0,1]$ and representing the frequency over time of an allele (or type) in a population of fixed size which evolves by resampling. We refer the reader to Etheridge's book~\cite{etheridge2011}. An important feature of these processes is that they model the phenomenon of random genetic drift, which is the fact that even in the absence of a selective advantage among the alleles (i.e. the model is neutral), allelic diversity will be reduced by the law of chance, so that the population will ultimately carry a single allele.

In this article, we consider Wright--Fisher processes with jumps in continuous time and space and generalise them by taking into account an extra force of selection.\linebreak Selection dynamics are typically modelled deterministically, so that the frequency of a type evolves both due to the resampling and due to a frequency-dependent term modeling how deleterious the allele considered is. Recently González and Spanò~\cite{zbMATH06873684} have established that discrete Wright--Fisher models with frequency-dependent selection can be rescaled to converge towards certain Markov processes called $\Xi$-Wright--Fisher process with frequency-dependent selection. In the latter work the approach of Neuhauser and Krone~\cite{zbMATH01078990} for modeling logistic selection is generalized by relating selection events with multiple (and not only binary) branching events in the ancestral genealogy. We shall focus on the simpler setting of $\Lambda$-Wright--Fisher process with frequency-dependent selection~($\Lambda$-WF processes with selection). In those processes, resampling events are simple in the sense that they only involve one fraction of the population.

Let $\Lambda$ be a finite measure over $[0,1]$. Let $\mu$ be a finite measure on $\bbN:=\{1,2,\,\ldots\}$. Denote by $f$ the generating function of the probability measure $\xi(\cdot)=\mu(\cdot)/\mu(\bbN)$ over $\bbN$, for all $x\in [0,1]$, $f(x):=\sum_{k=1}^{\infty}x^{k}\xi(k)$ and set $\sigma=\mu(\bbN)$. Consider the stochastic equation
\begin{multline}\label{SDEselecintro}
X_t(x)=x+\int_{0}^{t}\int_{0}^{1}\int_{0}^{1}z\left(\bbmun_{\{v\,\leq\,X_{s-}(x)\}}-X_{s-}(x)\right)\bar{\clM}(\ddr s,\ddr v, \ddr z)\\
-\sigma \int_{0}^{t}X_{s}(x)\big(1-f(X_{s}(x))\big)\ddr s,
\end{multline}
where $\clM$ is a Poisson point process on $\bbR_+\times [0,1]\times [0,1]$ with intensity $m(\ddr t,\ddr v, \ddr z)=\ddr t\otimes \ddr v\otimes z^{-2}\Lambda(\ddr z)$ and $\bar{\clM}$ stands for the compensated measure $\bar{\clM}=\clM-m$. Notice that the integrand in the stochastic integral with respect to $\bar{\clM}$ vanishes when the process reaches $0$ or $1$. Moreover, $f(0)=0$, $f(1)=1$ and for all $x\in [0,1]$, $1-f(x)\geq 0$ so that the drift term in~\eqref{SDEselecintro} is negative. In the general case, $\Lambda$-Wright--Fisher processes may have a diffusion part. We focus in this work on the case of a measure $\Lambda$ on $[0,1]$ with no mass at $0$ or at $1$.

Any process $(X_t(x),t\geq 0)$ solution to the equation~\eqref{SDEselecintro} is valued in $[0,1]$. Imagine a population of constant size $1$, whose individuals carry at any time one allele among a set of two alleles $\{a,A\}$. Suppose that the process $(X_t(x),t\geq 0)$ follows the frequency of allele $a$ when initially the proportion of individuals carrying allele $a$ is of size $x$. Before reaching boundaries, the time-dynamics of $(X_t(x),t\geq 0)$ consists of two parts:
\begin{itemize}
\item the resampling which is governed by the Poisson random measure $\clM$: for any $(t,v,z)$ atom of $\clM$,
\begin{itemize}
\item if $v\leq X_{t-}(x)$, then allele $a$ is sampled and a fraction $z\in (0,1)$ of the alleles $A$ at time $t-$ is replaced by the allele $a$ at time $t$. The frequency of allele $a$ increases:
\[
X_t(x)=z\big(1-X_{t-}(x)\big)+X_{t-}(x),
\]

\item if $v>X_{t-}(x)$, then allele $A$ is sampled and a fraction $z\in (0,1)$ of the alleles $a$ at time $t-$ is replaced by the allele $A$ at time $t$. The frequency of allele $a$ decreases:
\[
X_t(x)=(1-z)X_{t-}(x),
\]
\end{itemize}
\item the selection which is modeled by function $f$ characterizing the disadvantage of allele $a$: the frequency of allele $a$ decreases continuously in time along the negative deterministic drift:
\[
-\sigma X_t(x)\big(1-f(X_{t}(x))\big)\ddr t.
\]
\end{itemize}


When $\sigma=0$, the drift term in~\eqref{SDEselecintro} governing the selection disappears and the solution of~\eqref{SDEselecintro} becomes the classical $\Lambda$-Wright--Fisher process, see Bertoin and Le Gall~\cite{LGB2} and Dawson and Li~\cite{DawsonLi}, which represents the evolution of the frequency of a \emph{neutral} allele (or type) in a two-allele model evolving by resampling. In particular, when there is no selection term, the SDE~\eqref{SDEselecintro} has a pathwise unique strong solution and the boundaries $0$ and $1$ are both absorbing whenever they are reached. The event of absorption at $1$ (respectively at $0$) is called \emph{fixation} of the allele $a$ (respectively $A$) in the genetics terminology. It corresponds to the fact that all individuals have a common allele from a finite time almost surely. Bertoin and Le Gall~\cite{LGB2} have established that in the setting with no selection, fixation at one of the boundaries occurs almost surely if and only if the measure $\Lambda$ satisfies the following condition
\begin{equation}\label{cdipsi}
\sum_{n=2}^{\infty}\frac{1}{\Phi(n)}<\infty,
\end{equation}
where for any $n\geq 2$,
\begin{equation}\label{phi2}
\Phi(n):=\int_{(0,1)}\big((1-x)^{n}+nx-1\big)x^{-2}\Lambda(\ddr x).
\end{equation}
Condition~\eqref{cdipsi} is perhaps better known in the coalescent framework, as a necessary and sufficient condition for the $\Lambda$-coalescent to come down from infinity. A well-known cornerstone result in the coalescent theory states that any $\Lambda$-Wright-Fisher process satisfies a certain \emph{duality relationship} with a $\Lambda$-coalescent process, see Donnelly and Kurtz~\cite{DonnKurtz} and Bertoin and Le Gall~\cite{MR1990057}. Backgrounds on those results are given in Section~\ref{background}. We call $(N^{(n)}_t,t\geq 0,n\in \bbN)$ the block counting process of a $\Lambda$-coalescent started from $n$ blocks. For all $x\in [0,1]$ and $n\in \bbN$,
\begin{equation}\label{predualintro}
\bbE\left[X_t(x)^n\right]=\bbE\left[x^{N_t^{(n)}}\right].
\end{equation}
By letting $n$ go to $\infty$ in the identity~\eqref{predualintro}, we see that fixation at $1$ occurs if and only if the $\Lambda$-coalescent \emph{comes down from infinity}, in the sense that although it starts from infinitely many blocks, the number of blocks is finite at any strictly positive time.

One of the first models generalizing the $\Lambda$-Wright--Fisher process by incorporating selection is perhaps the logistic case for which $f(x)=x$ and the drift term in the SDE~\eqref{SDEselecintro} takes the form $-\sigma x(1-x)$. In this setting the measure $\mu$ reduces to a Dirac mass at $1$ with weight $\sigma$. Such processes have been studied by Baake et al.~\cite{baake2016}, Bah and Pardoux~\cite{zbMATH06403252}, Etheridge and Griffiths~\cite{zbMATH05832932}, Griffiths~\cite{zbMATH06387922} and Foucart~\cite{zbMATH06346921}. Bah and Pardoux~\cite[Theorem~4.3]{zbMATH06403252} have established that in the logistic case, fixation at $0$ or $1$ occurs almost surely if and only if~\eqref{cdipsi} is satisfied. In particular, when~\eqref{cdipsi} holds, despite that allele $a$ is deleterious when $\sigma>0$, the population still has a positive probability to get fixed on allele $a$ in a finite time almost surely.

The behavior of the positive function $x\mapsto 1-f(x)$ near $1$ actually reflects the strength of the selective advantage of allele $A$ over $a$. The question addressed in the present article is to see whether a selection term can overcome the $\Lambda$-resampling mechanism and prevent fixation of the deleterious allele $a$.

When $f$ is Lipschitz on $[0,1]$, i.e. $f'(1-)<\infty$, fundamental results on SDEs with jumps, see e.g.~\cite{DawsonLi}, entail that there exists a unique strong solution to~\eqref{SDEselecintro}. Moreover pathwise uniqueness holds and since $1$ is always a solution, it entails that the process is absorbed at $1$ if it reaches it. We shall actually see that in this case, fixation at boundary $1$ is always possible when~\eqref{cdipsi} holds.

When the drift term in~\eqref{SDEselecintro} is non-Lipschitz at $1$, namely $f'(1-)=\infty$, pathwise uniqueness of the solution to the SDE~\eqref{SDEselecintro} might not hold. In this case the only solution to~\eqref{SDEselecintro} whose existence and uniqueness is guaranteed is the minimal one, $(X_t^{\min},t\geq 0)$, which is stopped upon reaching boundary $1$. Several weak solutions to~\eqref{SDEselecintro} with different behaviors at boundary $1$ may exist. Let $(X^{\rmr}_t(x),t\geq 0,x\in [0,1])$ be a process valued in $[0,1]$. For any $x\in [0,1]$, let $\tau_1$ be the first hitting time of boundary $1$, i.e. $\tau_1:=\inf\{t>0: X^{\rmr}_t(x)=1\}\in [0,\infty]$. The process $(X^{\rmr}_t,t\geq 0)$ is said to be an extension of the minimal process $(X^{\min}_t,t\geq 0)$, if $(X^{\rmr}_{t\,\wedge\,\tau_1},t\geq 0)$ has the same law as $(X^{\min}_t,t\geq 0)$.

Following Feller's terminology for diffusions, see e.g. Karlin and Taylor's book~\cite[Chapter~15, Section~6]{MR611513}, a boundary is said to be natural if the process can not reach the boundary and can not leave it. The boundary is an exit if the process can reach the boundary but can not leave it. Symmetrically, it is said to be an entrance if the process can not access the boundary but leaves it; and finally the boundary is regular if the process enters into it and is able to get out from it.

In the sequel, we say that a boundary is \emph{absorbing} if when started from the boundary, the process stays at the boundary at any future time almost surely.\linebreak So that an exit boundary is always absorbing and a regular boundary is absorbing if it is subject to the prescription that the process fixes at the boundary once it is attained. In other words, the process with a \emph{regular absorbing} boundary is stopped at the boundary. By definition, an entrance boundary is non-absorbing as well as a regular boundary when the process is not stopped at it.

We will find new phenomena occurring in the presence of certain strong selection. In particular, despite the strength of the resampling rule under the condition~\eqref{cdipsi}, we shall find regimes for which even though the population starts entirely with the deleterious allele, its frequency will vanish in finite time almost surely. Namely, we will construct an extension of the minimal process with boundary $1$ as an entrance. Similarly, we will find sufficient condition under which the selection advantage for allele $A$ is not strong enough and for which boundary $1$ is an exit. Last but not least, we shall find regimes in which all individuals carry the deleterious allele for a set of times of negligible Lebesgue measure, before the selection starts to act effectively and that the deleterious allele vanishes.



Our method relies on the study of an extension constructed in the following way. We first look at processes, solution to the Equation~\eqref{SDEselecintro} with an additional drift term $-\lambda X_t\ddr t$ with $\lambda>0$. This drift can be seen as modeling \emph{mutation} from the deleterious allele $a$ to the advantaged one $A$. We shall see that under the assumption that there is no Kingman component, i.e. $\Lambda(\{0\})=0$, those processes, called $X^{\lambda}$'s, can all be started from boundary $1$. Our core object of study is the limit process that arises when the parameter $\lambda$ tends to $0$ (i.e. the mutation rate becomes very low). Hence define formally the limit process $X^{\rmr}$ as
\[
X^{\rmr}_t:=\underset{\lambda\,\rightarrow\,0+}{\lim} X^{\lambda}_t \text{ for all } t\geq 0.
\]
The convergence will be made precise later and we shall see that the process $X^{\rmr}$ is extending the minimal solution to~\eqref{SDEselecintro}. The possible behaviors at boundary $1$ of the extended process $(X_t^{\rmr},t\geq 0)$ are defined rigorously as follows.
Recall $\tau_1$ the first hitting time of boundary $1$. The $\Lambda$-WF process with selection has boundary $1$ accessible if for any $x\in (0,1)$, $\bbP_x(\tau_1<\infty)>0$. Furthermore, the boundary $1$ is
\begin{itemize}
\item[-] exit when $1$ is accessible and a.s. $X^{\rmr}_t(1)=1$ for all $t\geq 0$;
\item[-] entrance when $1$ is not accessible and a.s. $X^{\rmr}_t(1)<1$ for some $t>0$;
\item[-] regular non-absorbing when $1$ is accessible and a.s. $X^{\rmr}_{t}(1)<1\text{ for some }t>0$;
\item[-] natural when $1$ is not accessible and a.s. $X_t^{\rmr}(1)=1$ for all $t\geq 0$.
\end{itemize}

In the same fashion as in the case without selection, in which fixation at $1$ is linked to the coming down from infinity of the $\Lambda$-coalescent (i.e. $\infty$ is an entrance boundary), we will relate each boundary behavior of $X^{\rmr}$ to the boundary behavior of another dual process $(N^{(n)}_t,t\geq 0,n\in \bbN)$ with values in $\bar{\bbN}:=\bbN\cup \{\infty\}$. Loosely speaking, the process $(N^{(n)}_t,t\geq 0,n\in \bbN)$ can be seen as the functional of the block counting process of a \emph{simple} exchangeable fragmentation-coalescence (EFC) process with coalescence measure $\Lambda$ and splitting measure $\mu$. A simple EFC process $(\Pi(t),t\geq 0)$ is a partition-valued process in which coalescence occurs as in a $\Lambda$-coalescent and fragmentation dislocates a block chosen uniformly among all present blocks into $k$ sub-blocks at rate $\mu(k)$. The process $(N^{(n)}_t,t\geq 0,n\in \bbN)$ is thus generalizing the block counting process of a $\Lambda$-coalescent by allowing positive jumps from $n$ to $n+k$ at linear rate $n\mu(k)$ for any $k\in \bbN$. More backgrounds on EFC processes are provided in Section~\ref{backgroundEFC}.



Let $\zeta_\infty:=\inf\{t>0: N^{(n)}_{t-} \text{ or } N^{(n)}_{t}=\infty\}$ be the first explosion time of the block counting process. We define the stopped process 
\[
\left(N^{\min,(n)}_t,t\geq 0\right):= \left(N^{(n)}_{t\,\wedge\,\zeta_\infty},t\geq 0\right)
\]
and refer to it as the minimal block counting process. We will establish the following two duality relationships. For any $t\geq 0$, $x\in (0,1)$ and $n\in \bbN$,
\begin{equation}\label{dualityintro}
\bbE\left[\left(X^{\min}_t(x)\right)^n\right]=\bbE\left[x^{N^{(n)}_t}\right] \text{ and } \bbE\left[\left(X^{\rmr}_t(x)\right)^n\right]=\bbE\left[x^{N^{\min,(n)}_t}\right].
\end{equation}

The identities in~\eqref{dualityintro} can be seen as generalisations of the duality~\eqref{predualintro} to any generating function $f$ (including thus those with $f'(1-)=\infty$). They are stated in forthcoming Theorem~\ref{thmmomentdual1} and Theorem~\ref{thmmomentduality2} and have numerous consequences.

We will see for instance that if the boundary $\infty$ of the block-counting process $(N_t^{(n)},t\geq 0)$ is an entrance, an exit or is regular then the boundary $1$ for the extended process $(X_t^{\rmr},t\geq 0)$ is respectively an exit, an entrance or a regular boundary. We summarize the classification in Table~\ref{correspondance} in Section~\ref{mainresults}. We mention that such correspondences between entrance and exit boundaries for processes satisfying a duality relationship have been observed in other contexts, see the seminal work of Cox and R\"osler~\cite{zbMATH03828986}. Reminiscent classifications between boundaries have been established in Foucart~\cite{zbMATH07055671} for logistic continuous-state branching processes, see also Hermann and Pfaffelhuber~\cite{Hermann} and Berzunza Ojeda and Pardo~\cite{Berzunza}.

The behaviors of the block counting process of simple EFCs at boundary $\infty$ have been investigated in Foucart~\cite{cdiEFC} and Foucart and Zhou~\cite{explosion}. Sufficient conditions for the boundary $\infty$ to be entrance, exit or regular have been identified. By making use of the correspondences, we will be able to transfer the results on EFC processes to results for $\Lambda$-WF processes with selection and vice versa. Explicit conditions on the resampling measure $\Lambda$ and selection function $f$ are given for each possible behavior at boundary $1$. Those results are given in Theorems~\ref{suffcondexitf}, \ref{theorementrancef} and~\ref{regularcase}.

In a somehow parallel way of our study of the behaviors at the boundaries of the process $(X^{\rmr}_t,t\geq 0)$, we will see that fundamental properties of the block counting process of simple EFCs derive from these two duality relations. The first duality relation in~\eqref{dualityintro} allows one to establish the Markov and Feller properties of the process $ (N_t^{(n)}, t\geq 0)$ and to study its positive recurrence. These questions were not addressed in the previous works on EFCs, \cite{MR2110018,cdiEFC} and~\cite{explosion}, we refer the reader to~\cite[Remark~2.14]{cdiEFC}. Those results are stated in Theorem~\ref{Markovblockcounting}, Theorem~\ref{rec} and Corollary~\ref{inftyregularreflecting}.



The paper is organized as follows. The main results are stated in Section~\ref{mainresults}. Backgrounds on $\Lambda$-Wright--Fisher processes with and without selection are provided in Section~\ref{backgroundWF}. We verify in particular that the SDE~\eqref{SDEselecintro} admits a unique solution up to the first hitting time of the boundaries, this is the so-called minimal $\Lambda$-WF processes with selection. The moment-duality between the $\Lambda$-Wright--Fisher process without selection and the process counting the number of blocks in a $\Lambda$-coalescent is recalled in this section. Consequences of this duality relationships for the boundaries of the $\Lambda$-Wright--Fisher process without selection and the $\Lambda$-coalescent are reviewed. This will serve us as a guide for establishing the corresponding duality when fragmentation is taken into account. In Section~\ref{backgroundEFC}, we briefly recall the notion of exchangeable fragmentation-coalescence processes and describe the process of its number of blocks. Results from~\cite{cdiEFC} are gathered. In Section~\ref{sec1}, we establish the first duality identity in~\eqref{dualityintro} and deduce from it some new results for the EFC processes. In Section~\ref{sec2}, we proceed to the construction of the extension $(X_t^{\rmr},t\geq 0)$ of the minimal process. We then provide a theoretical classification of its boundary behaviors. In Section~\ref{sec3}, we go further into the classification by finding new correspondences between a regular reflecting boundary and a boundary which is regular for itself. Lastly, it is shown in this section, that when boundary $1$ is non-absorbing for the extended process, then under the assumption~\eqref{cdipsi} the latter gets absorbed at $0$ (i.e. fixation of the advantageous allele occurs). In Section~\ref{application}, we apply the results on simple EFCs, recalled in Section~\ref{backgroundEFC}, to the $\Lambda$-Wright--Fisher processes with selection and provide explicit sufficient conditions for each boundary behavior.


\subsection*{Notation}
We denote by $C([0,1])$ the space of continuous functions on $[0,1]$, $C^{2}((0,1))$ and $C^2_c([0,1])$ are respectively the space of twice continuously differentiable functions on $(0,1)$ and the space of twice continuously differentiable functions whose derivatives have compact support included in $(0,1)$. The integrability of a function $g$ in a left-neighbourhood of $a\in (0,\infty]$ is denoted by $\int^{a-}g(x)\ddr x<\infty$. We set $\bar{\bbN}:=\{1,2,\,\ldots,\,\infty\}$, the one-point compactification of $\bbN$, the latter is a compact set for the metric $d(n,m):=|n-m|$ for $n,m\in \bbN$ and $d(\infty,n):=1/n$ for any $n\in \bar{\bbN}$, where by convention $1/\infty =0$. The space of continuous functions on $\bar{\bbN}$ is denoted by $C(\bar{\bbN})$. A function $f$ belongs to $C(\bar{\bbN})$ if and only if $f(n)\underset{n\,\rightarrow\,\infty}{\longrightarrow} f(\infty)$. We denote by $\bbP_z$ the law of the process under consideration started from $z$, the corresponding expectation is $\bbE_z$.



\section{Main results}\label{mainresults}
Recall that $(X_t^{\min},t\geq 0)$ stands for the minimal $\Lambda$-Wright--Fisher process with frequency-dependent selection, that is to say the unique solution to~\eqref{SDEselecintro} which is absorbed at its boundaries after it reaches them. Existence and uniqueness of this process will be verified in Section~\ref{backgroundWF}, see Lemma~\ref{minimalprocess}. Recall that we denote by $(N_t^{(n)},t\geq 0)$ the block-counting process of a simple EFC started from $n$ blocks with coalescence measure $\Lambda$ and splitting measure $\mu$. We refer to Section~\ref{backgroundEFC} for more details.

We first state a duality relationship which holds for any minimal $\Lambda$-Wright--Fisher process with frequency-dependent selection, subject to the condition $\Lambda(\{1\})=0$. No assumption on the generating function $f$ is made.

\begin{theo}\label{thmmomentdual1}
The Markov process $(X_t^{\min}(x),t\geq 0,x\in [0,1])$ satisfies the following property. For any $t\geq 0,\, x\in [0,1]$ and any $n\in \bbN$
\begin{equation}\label{dualEFC}
\bbE\left[X^{\min}_t(x)^{n}\right]=\bbE\left[x^{N^{(n)}_t}\right].
\end{equation}
In particular,
\[
\bbP\left(X^{\min}_t(x)=1\right)=\underset{n\,\rightarrow\, \infty}{\lim}\bbE\left[x^{N^{(n)}_t}\right]= \bbE\left[x^{N^{(\infty)}_t}\right]\in [0,1]
\]
and the process $(X^{\min}_t(x),t\geq 0, x\in[0,1])$ gets absorbed at $1$ with positive probability if and only if the process $(N^{(\infty)}_t,t\geq 0)$ comes down from infinity (i.e. $\infty$ is either an entrance or a regular non-absorbing boundary).
\end{theo}

\begin{rema}
We shall see along the proof of Theorem~\ref{thmmomentdual1}, that~\eqref{dualEFC} can be generalized to cases where $\Lambda$ has a mass at $0$, see the forthcoming Remark~\ref{remKing}. We shall however focus in the sequel on pure-jump $\Lambda$-Wright--Fisher processes, namely those with $\Lambda(\{0\})=0$, see Remark~\ref{remwhynotkingman}.
\end{rema}

As explained in the introduction, we deduce from Theorem~\ref{thmmomentdual1} two important results for the block counting process of a simple EFC process. Those results were left unaddressed in~\cite{cdiEFC} and~\cite{explosion}.

\begin{theo}[Markov property of $(\#\Pi(t),t\geq 0)$]\label{Markovblockcounting}
Let $(\Pi(t),t\geq 0)$ be a simple EFC process whose coalescence measure is $\Lambda$ and splitting measure is $\mu$. The block counting process $(N_t,t\geq 0):=(\#\Pi(t),t\geq 0)$ with state-space $\bar{\bbN}$ is a Markov process satisfying the Feller property.
\end{theo}

Recall the first hitting times of the boundary $\tau_i:=\inf\{t\geq 0; X_t^{\min}=i\}$ for $i\in \{0,1\}$.
\begin{theo}[Recurrence of $(\#\Pi(t),t\geq 0)$] \label{rec}
If the process $(\#\Pi(t),\linebreak t\geq 0)$ comes down from infinity, then it is positive recurrent and has a stationary distribution whose generating function is $\varphi:x\in [0,1] \mapsto \bbP_x(\tau_1<\tau_0)$.
\end{theo}

\begin{rema}
A sufficient condition for coming down from infinity of simple EFC processes is given in~\cite[Theorem~1.1]{cdiEFC}. In the notation of~\cite{cdiEFC}, if $\theta^{\star}<1$, then the block counting process comes down from infinity and by Theorem~\ref{rec}, $(N_t,t\geq 0)$ is positive recurrent.
\end{rema}

\begin{rema}
The stationary distribution of $(N_t,t\geq 0)$ is carried over $\bbN$ if and only if the generating function $\varphi$ is non-defective i.e. $\varphi(x)=\bbP_x(\tau_1<\tau_0)\underset{x\rightarrow 1-}{\longrightarrow} 1$.
\end{rema}

The next theorem introduces the extended process $(X_t^{\rmr},t\geq 0)$. As explained in the introduction, it will be built from simpler processes $(X_t^{\lambda},t\geq 0)$ whose boundary $1$ is always an entrance. This limiting procedure is explained in the proof, see Lemma~\ref{recextensX}, which is deferred to Section~\ref{sec2}. As mentioned before, by regular absorbing boundary we mean that the process considered has its boundary regular but is stopped at it.

\begin{theo}\label{thmmomentduality2}
There exists a Feller process $(X_t^{\rmr}(x),t\geq 0, x\in [0,1])$ extending the minimal process such that for any $t\geq 0$ and any $n\in \bbN$,
\begin{equation}\label{dualityEFCref}
\begin{split}
\bbE\left[X^{\rmr}_t(x)^n\right] &=\bbE\left[x^{N^{\min,(n)}_t}\right] \text{ for any } x\in [0,1) \\
\text{ and}\quad \bbE\left[X^{\rmr}_t(1)^n\right] &=\bbP\left(N^{\min,(n)}_t<\infty\right),
\end{split}
\end{equation}
where $(N^{\min,(n)}_t,t\geq 0):=(N_{t\,\wedge\, \zeta_{\infty}}^{(n)},t\geq 0)$ for $\zeta_{\infty}:=\inf\{t>0: N_t^{(n)}=\infty\}\in [0,\infty]$. Moreover,
\begin{enumerate}\romanenumi
\item\label{theo2.7.i} the boundary $1$ is an entrance for $(X^{\rmr}_t(x),t\geq 0)$ if and only if $\infty$ is an exit for $(N^{(n)}_t,t\geq 0)$;
\item\label{theo2.7.ii} the boundary $1$ is regular non-absorbing for $(X^{\rmr}_t(x),t\geq 0)$ if and only if $\infty$ is regular absorbing for $(N^{\min,(n)}_t,t\geq 0)$, i.e. $\infty$ is regular non-absorbing for the process $(N^{(n)}_t,t\geq 0)$;
\item\label{theo2.7.iii} the boundary $1$ is an exit for $(X^{r}_t(x),t\geq 0)$ if and only if $\infty$ is an entrance of $(N^{(n)}_t,t\geq 0)$;
\item\label{theo2.7.iv} the boundary $1$ is natural for $(X^{r}_t(x),t\geq 0)$ if and only if $\infty$ is natural for $(N^{(n)}_t,t\geq 0)$.
\end{enumerate}
\end{theo}

The equivalences stated above will be established by making use of the two duality relationships~\eqref{dualEFC} and~\eqref{dualityEFCref}. The one-to-one correspondences between behaviors at the boundaries are summarized in the following table.


\begin{table}[htpb]
\caption{Classification of boundaries.}\label{correspondance}
\begin{center}
\begin{tabular}{|c|c|}
\hline Boundary $1$ of $X^{\rmr}$ & Boundary $\infty$ of $N$\\
\hline entrance & exit \\
\hline regular & regular \\
\hline exit & entrance \\
\hline natural & natural\\
\hline
\end{tabular}
\vspace*{2mm}
\end{center}
\end{table}

We stress that in the regular case, see the second line of Table~\ref{correspondance} above, if one process has its boundary regular \emph{non-absorbing} then according to Theorem~\ref{thmmomentduality2}-\eqref{theo2.7.ii}, its dual process has necessarily its boundary regular \emph{absorbing}.

Since there are several possible ways to leave a regular boundary, see e.g.~\cite[Chapter~15, Section~8]{MR611513} in the case of diffusions processes, Table~\ref{correspondance} does not specify completely the behavior of the process at the boundary when it is regular non-absorbing. Recall that a regular boundary is said to be \emph{reflecting} when the set of times at which the process lies at the boundary, has a zero Lebesgue measure.\linebreak A regular boundary is also said to be \emph{regular for itself} if the process started from the boundary returns immediately to it almost surely. In the same fashion, we classify the exit and entrance boundaries by saying that boundary $1$ is an \emph{instantaneous entrance} if it is an entrance and the first entrance time in $[0,1)$, $\tau^{1}:=\inf\{t>0: X^{\rmr}_t(x)<1\}$ satisfies $\bbP_1(\tau^{1}=0)=1$. The boundary $\infty$ is an \emph{instantaneous exit} if it is an exit and the first explosion time $\zeta_\infty:=\inf\{t>0: N^{(n)}_{t-} \text{ or } N^{(n)}_{t}=\infty\}$ satisfies for any $t>0$, $\bbP_n(\zeta_\infty \leq t)\longrightarrow 1$, as $n$ goes to $\infty$. Similar definitions hold for instantaneous exit boundary $1$ and instantaneous entrance boundary $\infty$.

The next theorem explains the possible behaviors of the dual processes at their boundaries when they are regular non-absorbing.

\begin{theo}[regular reflecting/regular for itself]\label{regularforitself}
The boundary $1$ of the extended process $(X_t^{\rmr},t\geq 0)$ is regular for itself (respectively, regular reflecting) if and only if the boundary $\infty$ of the process $(N_t,t\geq 0) $ is regular reflecting (respectively, regular for itself).
\end{theo}

\begin{prop}[Instantaneous entrance/exit]\label{instantanprop}
Assume that the boundary $1$ is an entrance for $(X_t^{\rmr},t\geq 0)$. The boundary $1$ is an instantaneous entrance if and only if $\infty$ is an instantaneous exit. Similarly, the boundary $1$ is an instantaneous exit if and only if $\infty$ is an instantaneous entrance.
\end{prop}

The next table summarizes Theorem~\ref{regularforitself} and Proposition~\ref{instantanprop}.
\begin{table}[htpb]
\caption{regular for itself/regular reflecting}\label{correspondancereg}
\begin{center}
\begin{tabular}{|c|c|}
\hline Boundary $1$ of $X^{\rmr}$ & Boundary $\infty$ of $N$ \\
\hline regular reflecting & regular for itself\\
\hline regular for itself & regular reflecting \\
\hline instantaneous entrance & instantaneous exit \\
\hline instantaneous exit & instantaneous entrance \\
\hline
\end{tabular}
\vspace*{2mm}
\end{center}
\end{table}

We address now the long-term behavior of the extended $\Lambda$-Wright--Fisher process with selection $(X^{\rmr}_t,t\geq 0)$. Recall that by fixation, we mean that all individuals get one of the two alleles and keep it forever. When $1$ is an exit, fixation of the deleterious allele has a positive probability to occur. When the boundary $1$ is regular non-absorbing or entrance, fixation at $1$ can not occur, and we shall actually see that there is almost sure fixation of the advantageous allele.

\begin{theo}\label{fixationat0}
Assume that $\Lambda$ satisfies~\eqref{cdipsi}. If $(X_t^{\rmr},t\geq 0)$
has boundary $1$ either regular non-absorbing or an entrance, then for all $x\in [0,1]$,
\[
\exists \ t_0>0; X_t^{\rmr}(x)=0 \text{ for all } t\geq t_0, \text{ a.s.}
\]
\end{theo}

Table~\ref{correspondance} and Table~\ref{correspondancereg} provide a theoretical classification of the boundaries. When there is selection, no necessary and sufficient conditions entailing that boundary $1$ of $X$ or boundary $\infty$ of $N$ is of a given type are known. We now identify explicit sufficient conditions on the resampling measure $\Lambda$ and the selection function $f$ for each possible boundary behavior. Proofs of those theorems are in Section~\ref{application}. They are obtained via the correspondences stated in Table~\ref{correspondance} and Table~\ref{correspondancereg}, by transferring previous results on the boundary $\infty$ of the block counting process $N$, obtained in~\cite{explosion} and summarized in Section~\ref{backgroundEFC}, to results on the boundary $1$ of $X^{\rmr}$.

Recall the map $\Phi$ defined in~\eqref{phi2}. Set $\Phi(z):=\Phi(\lfloor z \rfloor)$ for any $z\geq 2$. We refer the reader to~\cite[Section~2.2]{cdiEFC} and~\cite{limic2015} for fundamental properties of $\Phi$. In particular, we recall that $\Phi(z)\underset{z\rightarrow \infty}{\sim} \Psi(z)$ with $\Psi(z)=\int_{(0,1)}(e^{-xz}-1+xz)x^{-2}\Lambda(\ddr x)$. Moreover if for some $\beta \in (0,1)$
\begin{equation}\label{regularlambda}
\Lambda(\ddr x)=h(x)\ddr x,\text{ for } x\in [0,x_0] \text{ with } h(x)x^{\beta}\underset{x\,\rightarrow\, 0+}{\longrightarrow}\rho,
\end{equation}
then $\Phi(z)\underset{z\rightarrow \infty}{\sim} dz^{1+\beta}$ with $d=\rho\frac{\Gamma(1-\beta)}{\beta(\beta+1)}$, where $\Gamma$ is the Euler-Gamma function.

The next theorem provides conditions over the selection function $f$ and the resampling measure $\Lambda$ for $1$ to be an absorbing boundary for the (non-stopped) process $(X_t^{\rmr},t\geq 0)$, so that $1$ is either an exit or a natural boundary.

\begin{theo}\label{suffcondexitf}
If $f$ is Lipschitz on $[0,1]$, or if $x\mapsto 1-f(x)$ is regularly varying at $1$ with index $\alpha\in [0,1)$ and satisfies
\begin{equation}\label{integralform2}
\int^{1-}\frac{1-f(x)}{(1-x)^3\Phi\big(1/(1-x)\big)} \ddr x<\infty,
\end{equation}
then the boundary $1$ of $(X^{\rmr}_t,t\geq 0)$ is absorbing. Assume~\eqref{integralform2} holds true, then
\begin{enumerate}\romanenumi
\item\label{theo 2.11.i} if $\int^{\infty}\frac{\ddr z}{\Phi(z)}<\infty$, $1$ is an instantaneous exit boundary;
\item\label{theo 2.11.ii} if $\int^{\infty}\frac{\ddr z}{\Phi(z)}=\infty$, $1$ is a natural boundary.
\end{enumerate}
\end{theo}


We now provide a sufficient condition on the resampling measure $\Lambda$ and the selection function $f$ entailing that the process $(X_t^{\rmr},t\geq 0)$, has boundary $1$ as an entrance. Introduce the following condition over the function $x\mapsto 1-f(x)$:

\begin{enonce*}{Condition}
$\clH$: there exists a positive function $L$ defined on $(0,1)$ such that $\int^{1-} \frac{1}{L(x)}\ddr x<\infty$, the map $h:x\mapsto \frac{L(x)}{(1-x)\log\left(1/(1-x)\right)}$ is eventually non-decreasing in the neighbourhood of $1$ and
\begin{center} 
$1-f(x)\geq L(x)$ for $x$ close enough to $1$.
\end{center}
\end{enonce*}

\begin{rema}
Condition $\clH$ encompasses a regularity assumption on the difference quotient of the function $f$ near $1$. Indeed the condition on the map $h$ holds if the function $x\mapsto (1-f(x))/(1-x)\log(1/(1-x))$ stays above a non-decreasing function in some neighbourhood of $1$. In this case, Condition $\clH$ reduces to $\int^{1-}\frac{\ddr x}{1-f(x)}<\infty$.
\end{rema}

\begin{theo}\label{theorementrancef}
Assume Condition $\clH$ holds.
If
\begin{equation}\label{entrancecondu}
\frac{(1-x)^2\Phi\big(1/\log(1/x)\big)}{1-f(x)}\underset{x\,\rightarrow\,1-}{\longrightarrow} 0,
\end{equation}
then the boundary $1$ is an instantaneous entrance boundary.
\end{theo}


\begin{rema}
Since $\Phi$ is non-decreasing, $\Phi(1/\log(1/x))\leq \Phi(1/(1-x))$ for any $x\in [1/2,1)$.
\end{rema}

\begin{exam}
If there exist $c>0$ and $\alpha \in (0,1)$ such that $1-f(x)\geq c(1-x)^{\alpha}$ for $x$ close enough to $1$ then Condition $\clH$ is satisfied. If furthermore $\Phi\big(1/(1-x)\big)(1-x)^{2-\alpha}\underset{x\,\rightarrow\,1-}{\longrightarrow} 0$, then~\eqref{entrancecondu} holds. This is the case for instance when~\eqref{regularlambda} is fulfilled with $0<\beta<1-\alpha$ and $d>0$.
\end{exam}

\begin{theo}\label{regularcase}
Let $\alpha,\beta\in (0,1)$ and $\sigma,\rho>0$. Assume
\begin{equation}\label{regularlambdamu}
\Lambda(\ddr x)=h(x)\ddr x \text{ with } h(x)\underset{x\,\rightarrow\,0+}{\sim}\rho x^{-\beta} \text{ and } \sigma \big(1-f(x)\big)\underset{x\,\rightarrow\,1-}{\sim} \sigma (1-x)^{\alpha}.
\end{equation}
The boundary $1$ of $(X_t^{\rmr},t\geq 0)$ is classified as follows :
\begin{enumerate}\romanenumi
\item \label{theo2.17.i}if $\alpha+\beta<1$, then $1$ is an instantaneous entrance;
\vspace*{1mm}
\item\label{theo2.17.ii} if $\alpha+\beta>1$, then $1$ is an instantaneous exit;
\vspace*{1mm}
\item\label{theo2.17.iii} if $\alpha+\beta=1$ and further,
\vspace*{2mm}
\begin{itemize}
\item if $\sigma/\rho>\frac{\pi}{(2-\alpha)\sin(\pi \alpha)}$, then $1$ is an instantaneous entrance;
\vspace*{2mm}
\item if $\frac{1}{(1-\alpha)(2-\alpha)}<\sigma/\rho<\frac{\pi}{(2-\alpha)\sin(\pi \alpha)}$, then $1$ is regular reflecting;
\vspace*{2mm}
\item if $\sigma/\rho<\frac{1}{(1-\alpha)(2-\alpha)}$, then $1$ is an instantaneous exit.
\end{itemize}
\end{enumerate}
\end{theo}

\begin{rema}
Cases~\eqref{theo2.17.i} and~\eqref{theo2.17.ii} are consequences of Theorem~\ref{theorementrancef} and Theorem~\ref{suffcondexitf}. Important examples for which the condition~\eqref{regularlambdamu} hold are coalescence measures $\Lambda$ of the Beta form, $\Lambda(\ddr x)=\rho x^{-\beta}(1-x)^{a-1} \ddr x$ for $\beta\in (0,1)$ and $a>0$, and generating functions $f$ associated to Sibuya distribution, $f(x)=1-(1-x)^{\alpha}$ for $\alpha\in (0,1)$.
\end{rema}

\begin{theo}\label{regularforitselfstable}
Let $\sigma,\rho>0$ and $\alpha \in (0,1)$.
Assume 
\[
\Lambda(\ddr x)=h(x)\ddr x \text{ with } h(x)\underset{x\rightarrow 0+}{\sim}\rho x^{-(1-\alpha)}
\quad
\text{and} 
\quad
\sigma\big(1-f(x)\big)\underset{x\,\rightarrow\,1-}{\sim} \sigma (1-x)^{\alpha}.
\]
If $\frac{1}{(1-\alpha)(2-\alpha)}<\sigma/\rho<\frac{\pi}{(2-\alpha)\sin(\pi \alpha)}$, then the extended process $(X_t^{\rmr},t\geq 0)$ has its boundary $1$ regular for itself.
\end{theo}


\begin{rema}
Since the process $(X_t^{\rmr},t\geq 0)$ is Feller, when boundary $1$ is regular reflecting and regular for itself, standard theory, see e.g.~\cite[Chapter~IV]{Bertoin96} ensures the existence of a local time of the process $(X_t^{\rmr},t\geq 0)$ at $1$ whose inverse subordinator has no drift.
\end{rema}

By combining Theorem~\ref{regularforitselfstable} and Theorem~\ref{regularforitself}, we will obtain the following corollary for the block counting process $(N_t,t\geq 0)$ of a simple EFC process $(\Pi(t),t\geq 0)$ whose splitting measure $\mu$ and coalescence measure $\Lambda$ are regularly varying. Recall that by Theorem~\ref{Markovblockcounting}, $(N_t,t\geq 0):=(\#\Pi(t),t\geq 0)$ is a Markov process with state-space $\bar{\bbN}$. The first assertion~\eqref{coro2.21.i} below specifies the behavior of $(N_t,t\geq 0)$ when its boundary $\infty$ is regular and answers a question raised but left unaddressed in~\cite{explosion}.

\begin{coro}\label{inftyregularreflecting}
Let $\alpha\in (0,1)$. Assume $\Phi(n)\underset{n\,\rightarrow\,\infty}{\sim} dn^{2-\alpha}$ with $d>0$ and $\mu(n)\underset{n\,\rightarrow\,\infty}{\sim} \frac{b}{n^{1+\alpha}}$ with $b>0$.
\begin{enumerate}\romanenumi
\item\label{coro2.21.i} If $\frac{\alpha \sin(\pi \alpha)}{\pi}<b/d<\alpha(1-\alpha),$ then the boundary $\infty$ of the process $(\#\Pi(t),t\geq 0)$ is regular reflecting.
\item\label{coro2.21.ii} If $b/d<\alpha(1-\alpha)$, then the process $(\#\Pi(t),t\geq 0)$ is positive recurrent and admits a stationary distribution carried over $\bbN$.
\end{enumerate}
\end{coro}



\section{Background on \texorpdfstring{$\Lambda$}{Lambda}-WFs with selection and simple EFCs}\label{background}

\subsection{\texorpdfstring{$\Lambda$}{Lambda}-Wright--Fisher processes with selection and \texorpdfstring{$\Lambda$}{Lambda}-coalescent}\label{backgroundWF}

\subsubsection{A stochastic differential equation with jumps} 

We introduce the class of $\Lambda$-Wright--Fisher processes with frequency-dependent selection. As we shall use it later, we slightly generalize SDE~\eqref{SDEselecintro} by allowing the generating function driving selection to be defective, namely such that $f(1)<1$. Let $\mu$ be a finite measure over $\bar{\bbN}:=\{1,2,;\,\ldots,\,\infty\}$ and $\Lambda$ be a finite measure over $[0,1]$ with $\Lambda(\{0\})=\Lambda(\{1\})=0$. Let $f$ be the generating function of the probability law $\mu(\cdot)/\mu(\bar{\bbN})$ over $\bar{\bbN}$: for all $x\in [0,1]$, $f(x)=\sum_{k\,\in\,\bbN}x^{k}\mu(k)/\mu(\bar{\bbN})$. When the function $f$ is defective, $1-f(1)>0$ and this term corresponds to the mass at infinity for the probability distribution associated to $f$. Consider the stochastic equation
\begin{multline}\label{SDEselec}
X_t(x)=x+\int_{0}^{t}\int_{0}^{1}\int_{0}^{1}z\left(\bbmun_{\left\{v\,\leq\,X_{s-}(x)\right\}}-X_{s-}(x)\right)\bar{\clM}(\ddr s,\ddr v, \ddr z)\\
-\mu(\bar{\bbN})\int_{0}^{t}X_{s}(x)\big(1-f(X_{s}(x))\big)\ddr s.
\end{multline}
When there is no selection, i.e. $f(x)=1$ for all $x\in [0,1]$ and the drift term in Equation~\eqref{SDEselec} vanishes, the process $(X_t(x),t\geq 0)$ valued in $[0,1]$, is a martingale (this property in terms of the population model can be thought as the neutrality assumption between the two alleles) and has both boundaries $0$ and $1$ absorbing. Existence and weak uniqueness of the solution to the SDE~\eqref{SDEselec}, when $f\equiv 1$, has been established by Bertoin and Le Gall~\cite{LGB2} through a martingale problem. Set $q(v,x):=1_{\{v\,\leq\,x\}}-x$ for any $x\in [0,1]$. The generator of the $\Lambda$-Wright--Fisher process without selection is the operator $\clA$ defined as follows
\begin{align}\label{generator}
\clA g(x)&:=\int_{[0,1]\times [0,1]}\big(g(x+zq(v,x))-g(x)-zq(v,x)g'(x)\big) z^{-2}\Lambda(\ddr z)\ddr v.
\end{align}
Dawson and Li~\cite{DawsonLi} have studied the SDE~\eqref{SDEselec} through techniques different from Bertoin and Le Gall. Among other results, it is established that under some assumptions on the drift term, \eqref{SDEselec} admits a flow of strong solutions $(X_t(x),t\geq 0,x\geq 0)$ for which $x\mapsto X_t(x)$ is càdlàg.

Consider now the setting with selection. Any process solution to~\eqref{SDEselec} has a generator $\clA^{\rms}$ acting on $C^{2}((0,1))$ given by
\begin{equation}\label{generatorWFs}
\clA^{\rms}g(x):=\clA g(x)+\mu(\bar{\bbN})x(f(x)-1)g'(x) \text{ for any }x\in (0,1).
\end{equation}

\begin{lemm}\label{minimalprocess}
Let $f$ be any generating function (possibly defective). There exists a unique strong solution to~\eqref{SDEselec} up to the first hitting time of the boundaries.
\end{lemm}
\begin{proof}
For any $n\geq 1$, one can find a Lipschitz function $b^{n}$ on $[0,1]$ such that $b^{n}(x):=\mu(\bar{\bbN})x(f(x)-1)$ if $0\leq x\leq 1-1/n$ and $b^{n}(x):=0$ if $x\geq 1-1/2n$. Consider the stochastic equation
\begin{equation}\label{sdepre}
X_t(x)=x+\int_{0}^{t}\int_{0}^{1}\int_{0}^{1}z\left(\bbmun_{\left\{v\,\leq\,X_{s-}(x)\right\}}-X_{s-}(x)\right)\bar{\clM}(\ddr s,\ddr v, \ddr z)+\int_0^{t}b^{n}(X_s)\ddr s.
\end{equation}
We verify now that the latter equation has a pathwise unique strong solution by applying~\cite[Theorem~5.1]{li2012}. We follow the notation of~\cite{li2012} and check Conditions (3.a), (3.b) and (5.a). Set $U_0=[0,1]\times [0,1]$, for any $u\in U_0$, denote the coordinates of $u$ by $u=(v,z)$ and let $\mu_0(\ddr u)=\mu_0(\ddr v,\ddr z):=\ddr v \otimes z^{-2}\Lambda(\ddr z)$. Define $g_0(x,u)=z(\bbmun_{\{v\leq x\}}-x)$ and note that the stochastic equation above corresponds to the stochastic equation (2.1) in~\cite{li2012} with $\sigma\equiv 0, g_1\equiv 0$ and $b(x):=b^{n}(x)\leq 0$. Since $b^{n}$ is Lipschitz on $[0,1]$, condition (3.a) is satisfied (with $r_m(z)=r_1(z):=L_n z$ for any $z$ and any $m$, and where $L_n$ is the Lipschitz constant of $b_n$). Condition (3.b) is verified in~\cite[Corollary~6.2]{li2012}. It remains only to check Condition (5.a). For any $x\in \bbR$, $g_0(x,(u,z))^2\leq z^2$ and $b^{n}(x)^2\leq 1$, hence
\[
\int_0^1 g_0(x,(u,z))^2 \mu_0(\ddr u)+b^{n}(x)^2\leq \int_0^1 z^{2}z^{-2}\Lambda(\ddr z)+1= \Lambda((0,1))+1\leq K\left(1+x^2\right).
\]
Denote by $(X_t^{n},t\geq 0)$ the solution of~\eqref{sdepre}. Let $\tau_{1-1/n}:=\inf\{t>0: X_t^{n}> 1-1/n\}$. By pathwise uniqueness if $m<n$ then $X_t^{m}=X_t^{n}$ for $t\leq \tau_{1-1/m}\wedge \tau_0$. Note that $\tau_{1-1/m}\underset{m\,\rightarrow\,\infty}{\longrightarrow} \tau_1$; thus we can define a process $(X_t,t<\tau)$ such that $X_t=X_t^{n}$ for all $t\leq \tau_{1-1/n}$ and all $n\geq 1$. The process is solution to~\eqref{SDEselec} and the uniqueness of this minimal solution plainly holds.
\end{proof}

\begin{rema}\label{Lipschitzrem}
When the selection term satisfies the condition $f'(1-)<\infty$ (i.e. $\mu$ has a finite mean), the function $x\mapsto x(f(x)-1)$ is Lipschitz over $[0,1]$, and~\cite[Theorem~2.1]{DawsonLi} ensures that there is a unique pathwise strong solution to~\eqref{SDEselec}. It is worth noticing that in this case, since the constant process $1$ is solution to~\eqref{SDEselec}, pathwise uniqueness entails in particular that if the boundary $1$ is reached then the process is absorbed at $1$.
\end{rema}

We are interested in cases where $f$ is not globally Lipschitz on $[0,1]$, i.e. $f'(1-)=\infty$ and in possible extensions of the minimal process after time $\tau$. The minimal $\Lambda$-Wright--Fisher process with selection governed by $f$, denoted by $(X_t^{\min},t\geq 0)$, is the process solution to~\eqref{SDEselec} that is stopped at the boundary after time $\tau:=\inf\{t>0: X^{\min}_t \notin (0,1)\}$. A consequence of Lemma~\ref{minimalprocess} is that the minimal process, or minimal solution, is the unique solution to the following (stopped) martingale problem:
\[
(\MP): \forall g\in C^2_c([0,1]), \quad \left(g\left(X^{\min}_{t\,\wedge\,\tau}\right)-\int_{0}^{t}\clA^{\rms}g\left(X^{\min}_{s\,\wedge\,\tau}\right)\ddr s, t\geq 0\right) \text{ is a martingale,}
\]
where $X^{\min}_{t}=1$ if $t\geq \tau_1=\tau_1\wedge \tau_0$ and $X^{\min}_{t}=0$ if $t\geq \tau_0=\tau_1\wedge \tau_0$. By extension of the minimal process solution of~\eqref{SDEselec}, we mean here a process $(X_t,t\geq 0)$ with infinite life-time such that the process $(X_{t},t<\tau)$ has the same law as $(X^{\min}_{t},t<\tau)$.

An important tool used repeatedly in the proofs is a comparison theorem due to Dawson and Li for solutions of~\eqref{SDEselec}. Here, in their notation, $\sigma=g'_1=g''_1=0$ and the conditions (2-a)-(2-d) of~\cite[Theorem~2.2]{DawsonLi} are fulfilled.

\begin{theo}[Theorem~2.2 in~\cite{DawsonLi}]\label{comparisontheorem}
Given two Lipschitz generating functions $f_1$ and $f_2$ such that $f_1\leq f_2$ (thus, by the Lipschitz assumption $f'_i(1-)<\infty$ for $i=1,2$) and two initial values $x_1\leq x_2$ in $[0,1]$, if $(X^i_t(x_i),t\geq 0)$ is the solution of~\eqref{SDEselec} with $f=f_i$ for $i=1,2$ and initial value $x_i$, then almost surely for all $t\geq 0$, $X^{1}_t(x_1)\leq X^{2}_t(x_2)$.
\end{theo}


\subsubsection{\texorpdfstring{$\Lambda$}{Lambda}-Wright--Fisher processes without selection}

A fundamental property of $\Lambda$-Wright--Fisher processes without selection is their link with the processes called $\Lambda$-coalescents. Those processes are valued in the space of partitions of $\bbN$ and are evolving by multiple (not simultaneous) mergings of equivalence classes (called blocks). The $\Lambda$-coalescents can be thought as representing the genealogy backwards in time of the ancestral lineages in the population evolving by resampling, see~\cite{MR1990057}. More backgrounds about $\Lambda$-coalescents are provided in Section~\ref{backgroundEFC}. Let $(\Pi(t),t\geq 0)$ be a $\Lambda$-coalescent and $(\#\Pi(t),t\geq 0)$ be its block counting process. This is a Feller process valued in $\bar{\bbN}$ whose generator is $\clL^{c}$ with for any $g:\bbN\mapsto \bbR$
\begin{align}
\clL^c g(n){:=}&\sum_{k=2}^{n}\binom{n}{k}\lambda_{n,k}(g(n-k+1)-g(n)),\label{coalpart}
\\
\intertext{where}
\lambda_{n,k}{:=}&\int_0^1 z^{k}(1-z)^{n-k}z^{-2}\Lambda(\ddr z) \text{ for } 2\leq k\leq n.\label{lambdank}
\end{align}
Recall Condition~\eqref{cdipsi}. Schweinsberg~\cite{CDI} has established that~\eqref{cdipsi} is necessary and sufficient for $\infty$ to be an entrance boundary of the block counting process, in this case although $\#\Pi(0)=\infty$, $\#\Pi(t)<\infty$, for any $t>0$ almost surely.

The following identity, established for instance in~\cite[Theorem~1, Equation~(8)]{LGB2}, links the block counting process of a pure $\Lambda$-coalescent with the $\Lambda$-Wright--Fisher process without selection through a moment-duality relationship. For any $x\in [0,1]$ and $n\in \bar{\bbN}$, if $\#\Pi(0)=n$ then
\begin{equation}\label{dualclassic}
\bbE\left[X_t(x)^n\right]=\bbE_n\left[x^{\#\Pi(t)}\right],
\end{equation}
where we denote by $\bbE_n$ the expectation conditionally given that $\#\Pi(0)=n$.


The identity~\eqref{dualclassic} has many important consequences. Firstly, since $X_t(x)$ is a bounded random variable, its law is entirely characterized by its moments and therefore the one-dimensional laws of $(\#\Pi(t),t\geq 0)$ are in one-to-one correspondence with those of the process $(X_t(x),t\geq 0)$. From a theoretical point of view one can see~\eqref{dualclassic} as a representation of the semigroup of the process $(X_t(x),t\geq 0)$. Moreover, for any $n\in \bbN$, $\#\Pi(t)\leq \#\Pi(0)=n$ for all $t\geq 0$, $\bbP_n$-a.s. and letting $x$ approach $1$ provides the identity
\[
\underset{x\,\rightarrow\,1-}{\lim}\bbE\left[X_t(x)^n\right]=\underset{x\rightarrow 1-}{\lim}\bbE_n\left[x^{\#\Pi(t)}\right]=\bbP_n(\#\Pi(t)<\infty)=1.
\]
Since $x\mapsto X_t(x)$ admits left-limits, we see from the above convergence and Lebesgue's theorem that $X_t(1-)=1$ almost surely for all $t\geq 0$, so that the boundary $1$ is not an entrance of the $\Lambda$-Wright--Fisher process, but is absorbing whenever it is reached. Letting $n$ go to infinity in~\eqref{dualclassic} entails the identity
\begin{equation}\label{classicid}
\bbP_x(\tau_1\leq t)=\bbE_{\infty}\left[x^{\#\Pi(t)}\right],
\end{equation}
where $\tau_1:=\inf\{t>0: X_t(x)=1\}$ and $(\#\Pi(t),t\geq 0)$ starts from $\infty$ under $\bbP_\infty$. When Schweinsberg's condition~\eqref{cdipsi} holds, $\#\Pi(t)<\infty$ a.s. for any $t>0$ and~\eqref{classicid} ensures that $\tau_1<\infty$ with positive probability. The identity~\eqref{classicid} provides a representation of the cumulative distribution function of $\tau_1$ and in a dual way a representation of the entrance law at $\infty$ of the process $(\#\Pi(t),t\geq 0)$. Since $1$ is absorbing, the event $\{\tau_1<\infty\}$ coincides with the event of fixation at $1$: $\{\exists \ t_1\geq 0; \forall t\geq t_1; X_t(x)=1\}$.


\subsubsection{Pure selection process}

When there is no resampling in the population, the frequency of the deleterious allele $(X_t(x),t\geq 0)$ solves the ODE
\begin{equation}\label{ode}
X_{t}(x)=x-\mu(\bar{\bbN})\int_{0}^{t}X_s(x)\big(1-f(X_s(x))\big)\ddr s.
\end{equation}
Equivalently, the map $(X_t(x),t\geq 0)$ satisfies for all $x\in (0,1)$ and $t\geq 0$,
\begin{equation}\label{integralform}
\int_{X_t(x)}^{x}\frac{\ddr u}{u(1-f(u))}=\mu(\bar{\bbN})t.
\end{equation}
A study of~\eqref{integralform} when $x$ tends to $1$, yields the following dichotomy. Either $\int^{1-}\frac{\ddr u}{1-f(u)}\linebreak=\infty$, and in order for the integral in~\eqref{integralform} to retain the value $\mu(\bar{\bbN})t$, we must have $X_t(1):=\underset{x\,\rightarrow\,1-}{\lim}X_t(x)=1$, or else $\int^{1-}\frac{\ddr u}{1-f(u)}<\infty$ and we must have
\[
X_t(1):=\underset{x\,\rightarrow\,1-}{\lim}X_t(x)\in (0,1).
\]
In the latter case, the function $(X_t(1),t\geq 0)$ solves~\eqref{ode} and starts from $1$. Hence, if $f(1)=1$ then~\eqref{ode} has two distinct solutions started from $1$, $X_t:=1$ for all $t\geq 0$ and $X_t:=\underset{x\,\rightarrow\,1-}{\lim} X_t(x)$ for all $t\geq 0$, if and only if
\begin{equation}\label{dynkin}
\int^{1-} \frac{\ddr x}{1-f(x)}<\infty.
\end{equation}
The integral in~\eqref{dynkin} converges for example when
\[
1-f(x)\underset{x\,\rightarrow\,1-}{\sim} \sigma (1-x)^{\alpha}
\]
with $\alpha\in (0,1)$ and $\sigma>0$. Note that if $f(1)<1$, i.e. $\mu(\infty)>0$, then~\eqref{dynkin} clearly holds.

Solutions to~\eqref{ode} are well-known in the theory of branching processes. Consider a continuous-time discrete-state branching process $(N_t^{(n)},t\geq 0)$ started from\linebreak $n\in \bbN$ with offspring measure $\mu$. Namely, the process jumps from any integer state $n$ to $n+k$ at rate $n\mu(k)$ for all $k\in \bbN$. We refer for instance to Harris' book\linebreak\cite[Chapter~V]{Harris}. Note that $\mu$ gives no mass at $-1$, so that there is no death in the process $(N_t^{(n)},t\geq 0)$ and its sample paths are almost surely non decreasing. The branching property of the process $(N_t^{(n)},t\geq 0)$ ensures that the boundary $\infty$ is absorbing whenever it is reached. One can identify the function $(X_t(x),x \in [0,1])$ with the generating functions of $N^{(1)}_t$ at any time $t$, more precisely for any $x\in [0,1]$, any $n\in \bbN$ and any $t\geq 0$
\begin{equation}\label{dualbranching}
X_t(x)^{n}=\bbE\left[x^{N_t^{(n)}}\right].
\end{equation}
Letting $x$ go towards $1$ in the identity above yields: 
\begin{align*}
\bbP_n(\zeta_\infty\leq t)=& \underset{x\,\rightarrow \,1-}{\lim} X_t(x) \text{ for any }t\geq 0,\\
\intertext{where}
\zeta_\infty{:=}& \inf\left\{t>0: N_{t-}^{(n)} \text{ or } N_t^{(n)}=\infty\right\}.
\end{align*}
This is the standard method for studying the explosion of process $(N_t^{(n)},t\geq 0)$, see~\cite[Chapter~V, Theorem~9.1]{Harris}. In particular, the boundary $\infty$ is an exit for the branching process $(N_t^{(n)},t\geq 0)$ if and only Dynkin's condition~\eqref{dynkin} holds.

When resampling and selection are both taken into account, the process solution to~\eqref{SDEselec} may have more involved boundary conditions than what has just been seen for processes without selection or without resampling. In a similar fashion as for the pure selection process and for the pure resampling one, we will see that the moments of a $\Lambda$-WF process with selection $(X_t,t\geq 0)$ can be represented via certain continuous-time Markov chains, with values in $\bar{\bbN}$, whose jumps are mixture of those of the block counting process of a $\Lambda$-coalescent and those of an increasing branching processes. These processes appear when counting the number of blocks in certain exchangeable partition-valued processes.

\subsection{Backgrounds on exchangeable fragmentation coalescence (EFC) processes}\label{backgroundEFC}


EFC processes, introduced by Berestycki in~\cite{MR2110018}, are Markov processes with state-space $\clP_\infty$, the space of partitions of $\bbN$. The number of non-empty blocks (i.e. equivalence classes) of a partition $\pi\in \clP_\infty$ is denoted by $\#\pi$. By definition an EFC process $(\Pi(t),t\geq 0)$ satisfies the following conditions :
\begin{enumerate}\romanenumi
\item for any time $t\geq 0$, the random partition $\Pi(t)$ is exchangeable, i.e. its law is invariant by the action of permutations with finite support;
\item it evolves in time by merging of blocks or fragmentation of an individual block into sub-blocks.
\end{enumerate}
We consider here only the subclass of simple EFC processes in which as in a\linebreak $\Lambda$-coalescent, there is no simultaneous multiple mergings, fragmentations occur at \emph{finite rate} and fragmentate any blocks into sub-blocks of infinite size (no formation of singletons). We shall not introduce the whole framework of partition-valued\linebreak processes here, for which we refer to Berestycki~\cite{MR2110018}, but focus on the block counting process.

We briefly explain below the dynamics of the block counting process with the help of two Poisson point processes. We focus on EFC processes whose initial partitions have all blocks of infinite size. In particular at all time in the system, there is no singleton blocks. We refer to~\cite[Sections~2.1 and 2.2]{cdiEFC} for details.

Let $\Lambda$ be a finite measure on $(0,1)$ and $\mu_{\Frag}$ be a finite (exchangeable) measure on $\clP_\infty$. We call $\Lambda$ and $\mu_{\Frag}$ respectively the coalescence measure and the fragmentation measure. Consider two independent Poisson point processes $\PPP_C$ and $\PPP_F$ respectively on $\bbR_+\times [0,1]^{\bbN}$ and $\bbR_+\times \clP _\infty\times \bbN$.

The intensity of $\PPP_C$ is $\ddr t\otimes \Ber_{\Lambda}$ where $\Ber_{\Lambda}$ stands for the law of an infinite exchangeable sequence $(X_i,i\geq 1)$ mixture of i.i.d Bernoulli random variables whose parameter $\underset{n\,\rightarrow\,\infty}{\lim}\frac{1}{n}\sum_{i=1}^{n}X_i \text{ has for ``intensity'' the measure } x^{-2}\Lambda(\ddr x).$

The intensity of $\PPP_F$ is $\ddr t\otimes \mu_{\Frag}(\ddr \pi)\otimes \#$ where $\#$ denotes the counting measure on $\bbN$ and $\mu_{\Frag}$ is a finite measure on $\clP_\infty$, the so-called fragmentation measure. Let $\mu$ be the image of $\mu_{\Frag}$ by the map $\pi \mapsto \#\pi-1$. We call $\mu$ the splitting measure.

\paragraph{Coalescence} Upon the arrival of an atom $(t,(X_i)_{i\,\geq\,1})$ of $\PPP_C$, given $\#\Pi(t-)\linebreak=n$, all blocks whose index $j\in [n]$ satisfies $X_{j}=1$ are merged. Given the parameter $x$ of the $X_i$'s, the number of blocks that merge at time $t$ has a binomial law with parameters $(n,x)$. Therefore, for any $k\in [|2,n|]$ the jump $\#\Pi(t)=\#\Pi(t-)-(k-1)$ has rate $\binom{n}{k}\lambda_{n,k}$ where we recall $\lambda_{n,k}:=\int_{[0,1]}x^{k}(1-x)^{n-k}x^{-2}\Lambda(\ddr x)$.
\vspace*{1mm}

\paragraph{Fragmentation} Upon the arrival of an atom $(t,\pi^{f},j)$ of $\PPP_F$, given $\#\Pi(t-)=n$, and $j\leq n$, then the $j^{\mathrm{th}}$-block of $\Pi(t-)$ is fragmented into $\#\pi^{f}\in \bar{\bbN}$ sub-blocks. Therefore, at time $t$, $\#\Pi(t)=\#\Pi(t-)-1+\#\pi^{f}=\#\Pi(t-)+k$ with $k=\#\pi^{f}-1$. Since there are $n$ blocks at time $t-$, the total rate at which such a jump occurs is $n\mu(k)$ for any $k\in \bar{\bbN}$.

\begin{lemm}[Proposition~2.11 in~\cite{cdiEFC}]\label{genEFC}
If $\Pi(0)$ has blocks of infinite size, then the process $(\#\Pi(t),t\geq 0)$ is a right-continuous process valued in $\bar{\bbN}$ and has the Markov property when lying on $\bbN$: i.e. setting $\zeta_{\infty}:=\inf\{t>0: \#\Pi(t-)=\infty \text{ or } \#\Pi(t)=\infty\}$, $(\#\Pi(t),t<\zeta_{\infty})$ is a continuous-time Markov chain whose generator is $\clL=\clL^c+\clL^f$ with $\clL^{c}$ is given in~\eqref{coalpart} and $\clL^f$ is acting on any bounded function $g$, with a limit at $\infty$ when $\mu(\infty)>0$, by
\begin{equation}\label{fragpart}
\clL^f g(n):=\sum_{k=1}^{\infty}n\mu(k)(g(n+k)-g(n))+n\mu(\infty)(g(\infty)-g(n)).
\end{equation}
\end{lemm}



Any EFC process $(\Pi(t),t\geq 0)$ is defined on the whole half-line of time, so that the block-counting process is well-defined past explosion (if explosion occurs). Sufficient conditions for the process $(N_t,t\geq 0):=(\#\Pi(t),t\geq 0)$ to have boundary $\infty$ absorbing or not have been found in~\cite{cdiEFC}. More precisely, the process can have the boundary $\infty$ exit, entrance or even regular for certain heavy-tailed splitting measures. We mention that when the block counting process $(N_t,t\geq 0)$ comes down from infinity (i.e. when $\infty$ is an entrance or a regular boundary), then since by assumption $\Lambda(\{1\})=0$, the process leaves the boundary $\infty$ instantaneously, see~\cite[Lemma~2.5]{cdiEFC}.

\begin{lemm}[{\cite[Corollary~1.2-(2)]{cdiEFC}}]\label{key}
Let $\lambda>0$. If the coalescence measure has no Kingman part, i.e. $\Lambda(\{0\})=0$, and $\mu(\infty)=\lambda$, then the process $(N_t,t\geq 0):=(\#\Pi(t),t\geq 0)$ has boundary $\infty$ as an exit, that is to say, for any $t\geq \zeta_{\infty}$, $N_t=\infty$ a.s.
where $\zeta_{\infty}:=\inf\{t>0: N_t=\infty\}$.
\end{lemm}

\begin{rema}\label{remkey}
When $\Lambda(\{0\})=c_{\rmk}>0$, the process $(N_t,t\geq 0)$ comes down from infinity (i.e leaves $\infty$) if and only if $\frac{2\lambda}{c_{\rmk}}<1$, see Kyprianou et al. \cite[Theorem~1.1]{kyprianou2017} and~\cite[Corollary~1.2-(1)]{cdiEFC}. Lemma~\ref{key} is a key lemma in our construction of an extension getting out from $1$ of the minimal $\Lambda$-WF process with selection.
\end{rema}

The next lemmas provide sufficient condition for $\infty$ to be an entrance, an exit or a regular boundary. They will be used in Section~\ref{application} in a ``dual'' way for the $\Lambda$-WF processes with selection. Recall that $\bar\mu$ denotes the tail of the splitting measure $\mu$, for any $n\in \bbN$, $\bar\mu(n)=\mu(\{n,n+1,\cdots\})$.
\begin{lemm}[{\cite[Theorem~3.4]{explosion}}] \label{nonexplosioncriterion}
Assume $\mu(\infty)=0$. If
\[
\sum_{n=2}^{\infty}\frac{n}{\Phi(n)}\bar{\mu}(n)<\infty,
\]
then $\infty$ is inaccessible for the process $(N_t,t\geq 0)$. Moreover, in this case,
\begin{enumerate}\romanenumi
\item if $\sum_{n=2}^{\infty}\frac{1}{\Phi(n)}<\infty$ then $\infty$ is an entrance boundary,
\item if $\sum_{n=2}^{\infty}\frac{1}{\Phi(n)}=\infty$ then $\infty$ is a natural boundary.
\end{enumerate}
\end{lemm}

The next lemma provides a sufficient condition for $\infty$ to be an exit boundary when $\mu(\infty)=0$. We define $\ell :n\mapsto \sum_{k=1}^{n}\bar{\mu}(k)$ and set the following condition on $\ell$.

\begin{enonce*}{Condition $\bbH$} 
There exists an eventually non-decreasing positive function $g$ on $\bbR_+$ such that $\int^{\infty}\frac{1}{xg(x)}\ddr x<\infty$ and
\begin{equation*}
\ell(n)\geq g(\log n)\log n \text{ for large } n.
\end{equation*}
\end{enonce*}

When Condition $\bbH$ is in force, \cite[Theorem~3.1]{explosion} provides a sufficient condition for boundary $\infty$ to be an exit. For the sake of simplicity, we shall work with the following direct corollary of this theorem:

\begin{lemm}[{\cite[Theorem~3.1, case $\rho=0$]{explosion}}]\label{suffcondexitN}
Assume Condition $\bbH$ holds and $\mu(\infty)=0$. If
\begin{equation*}\label{suffcondentrancePhi}
\underset{n\,\rightarrow\,\infty}{\lim}\ \frac{\Phi(n)}{n\ell(n)}=0,
\end{equation*}
then the process $(N_t,t\geq 0)$ has $\infty$ as an exit boundary.
\end{lemm}

We now give the classification of the boundary $\infty$ in some regularly varying cases found in~\cite{explosion}.

\pagebreak
\begin{lemm}[{\cite[Theorem~3.7]{explosion}}]\label{stablefragtheorem}
Assume that $\Phi(n)\underset{n\,\rightarrow\,\infty}{\sim} dn^{1+\beta}$ with $d>0$ and $\beta\in (0,1)$ and $\mu(n)\underset{n\,\rightarrow\,\infty}{\sim} \frac{b}{n^{\alpha+1}}$ with $b>0$ and $\alpha\in (0,\infty)$. Then
\begin{enumerate}\romanenumi
\item\label{lemm3.9.I} if $\alpha+\beta<1$, then $\infty$ is an exit boundary,
\vspace*{1mm}
\item\label{lemm3.9.II} if $\alpha+\beta>1$, then $\infty$ is an entrance boundary,
\vspace*{1mm}
\item\label{lemm3.9.III} if $\alpha+\beta=1$ and further,
\vspace*{2mm}
\begin{itemize}
\item if $b/d>\alpha(1-\alpha)$, then $\infty$ is an exit boundary,
\vspace*{2mm}
\item if $\frac{\alpha \sin(\pi \alpha)}{\pi}<b/d<\alpha(1-\alpha)$, then $\infty$ is a regular boundary,
\vspace*{2mm}
\item if $b/d<\frac{\alpha \sin(\pi \alpha)}{\pi}$, then $\infty$ is an entrance boundary.
\end{itemize}
\end{enumerate}
\end{lemm}

Several coupling procedures have been designed in~\cite{cdiEFC} in order to study the process $(N_t,t\geq 0):=(\#\Pi(t),t\geq 0)$ started from $\infty$. At several places later, we will use a monotone coupling of the block counting process $(\#\Pi(t),t\geq 0)$ in the initial values.

Let $n\in \bar{\bbN}$ and consider the process $(\Pi^{(n)}(t),t\geq 0)$ which starts from the $n$ first blocks of $\Pi(0)$, i.e. $\Pi^{(n)}(0)=\{\Pi_1(0),\,\ldots,\,\Pi_n(0)\}$, and evolves along the same Poisson Point Processes $\PPP_C$ and $\PPP_F$ as $(\Pi(t),t\geq 0)$. We refer to~\cite[Lemma~3.3 and Lemma~3.4]{cdiEFC} for details on the construction. The process $(\Pi^{(n)}(t),t\geq 0)$ follows the fragmentations and coagulations in the system restricted to the integers belonging to $\cup_{i=1}^{n}\Pi_i(0)$. In the sequel, we write $(N_t^{(n)},t\geq 0):=(\#\Pi^{(n)}(t),t\geq 0)$ for the process counting the blocks of $(\Pi^{(n)}(t),t\geq 0)$. Notice that by definition, $(\Pi^{(\infty)}(t),t\geq 0)$ coincides with $(\Pi(t),t\geq 0)$ and thus $(N_t^{(\infty)},t\geq 0)=(N_t,t\geq 0)$.

The process $(N_t^{(n)},t\geq 0)$ is at the core of our study, and from now on we simply call it ``block counting process''.
\begin{lemm}[{Monotonicity in the initial values,\cite[Lemma~3.4]{cdiEFC}}]\label{monotonicity}
For any $n\geq 1$,
\begin{equation*}
N_t^{(n)}\leq N^{(n+1)}_t\text{ and }N_t^{(n)}\underset{n\,\rightarrow\,\infty}{\longrightarrow} N^{(\infty)}_t\text{ for all }t\geq 0\text{ a.s.}
\end{equation*}
Moreover, the process $(N^{(n)}_t,t\zeta_{\infty})$ is Markovian and has the same law as $(N_t,t<\zeta_{\infty})$ when $N_0=n<\infty$.
\end{lemm}

The next lemma ensures that one can approach from below the process $(N_t^{(n)},t\geq 0)$ by a non-decreasing sequence of non-explosive processes. For any $m\in \bbN$, set $\bar{\mu}(m):=\mu(\{m,\,\ldots,\,\infty\})$.

\begin{lemm}[{\cite[Lemma~3.8]{cdiEFC}}] \label{Nnm}
For any $n\in \bar{\bbN}$, there exists a non-decreasing sequence of processes $(N^{(n)}_m(t),t\geq 0)_{m\,\geq\,1}$ started from $n\in \bar{\bbN}$ such that
\begin{enumerate}\romanenumi
\item\label{lemm3.11.i} for any $m\geq 1$, $(N^{(n)}_m(t),t\geq 0)$ has generator, the operator $\clL^{m}$ acting on any function $g: \bbN\rightarrow \bbR_+$ as follows
\begin{equation}\label{genNm}
\clL^{m}g(\ell):=\clL^{c}g(\ell)+\ell\sum_{k=1}^{m}\mu_{m}(k)(g(\ell+k)-g(\ell)), \text{ for all } \ell \in \bbN,
\end{equation}
where $\mu_{m}(k):=\mu(k)$ if $k\leq m-1$ and $\mu_m(m):=\bar{\mu}(m)$;
\item\label{lemm3.11.ii} for any $m\in \bbN$, the process $(N_m^{(n)}(t),t\geq 0)$ does not explode;
\item\label{lemm3.11.iii} almost surely for any $n\in \bar{\bbN}$, $m\in \bbN$ and all $t\geq 0$, $N^{(n)}_{m}(t)\leq N^{(n+1)}_{m}(t)$,
and
\[
\underset{m\,\rightarrow\,\infty}{\lim} N^{(n)}_m(t)=N^{(n)}_t \text{a.s.}
\]
\end{enumerate}
\end{lemm}


\section{Proofs of Theorems~\ref{thmmomentdual1}, \ref{Markovblockcounting}, \ref{rec}}\label{sec1}



In this section we focus on the study of the minimal $\Lambda$-WF process with selection, $(X_t^{\min},t\geq 0)$. Lemma~\ref{minimalprocess} ensures its existence and uniqueness without assumption on the splitting measure $\mu$. We establish Theorem~\ref{thmmomentdual1} where a first duality relationship~\eqref{dualEFC} between $(X_t^{\min},t\geq 0)$ and the block counting process $(N_t,t\geq 0)$ is displayed. This will generalize the duality relationships~\eqref{dualclassic} and~\eqref{dualbranching} known for $\Lambda$-coalescents and branching processes. To prove Theorem~\ref{thmmomentdual1}, we shall need the monotone coupling $(N_t^{(n)},t\geq 0, n\in \bar{\bbN})$ and $(N_m^{(n)}(t),t\geq 0, n\in \bar{\bbN})$ described in Lemmas~\ref{monotonicity} and~\ref{Nnm}. Once Theorem~\ref{thmmomentdual1} established, consequences (Theorems~\ref{Markovblockcounting} and~\ref{rec}) for the block counting process $(N_t,t\geq 0)$ are deduced.



Recall $g_n(x)=g_x(n)=x^n$ for all $x\in [0,1]$ and $n\in \bbN$. We work with the convention
\begin{equation}\label{convention}
\underset{n\,\rightarrow\,\infty}{\lim} x^{n}=\bbmun_{\{x=1\}}, \text{ and } \underset{x\,\rightarrow\,1-}{\lim}x^{n}=1_{\{n\,<\,\infty\}}.
\end{equation}

\subsection{Proof of Theorem~\ref{thmmomentdual1}}\label{momentdual1}

\begin{proof}
Let $\Lambda$ be a coalescence measure such that $\Lambda(\{1\})=0$ and $\mu$ be a splitting measure with possibly a mass at $\infty$. For any $x\in [0,1]$, recall that $(X_t^{\min}(x),t\geq 0)$ is the minimal process which is absorbed at its boundary once it has reached it.



The following is the scheme of the proof of Theorem~\ref{thmmomentdual1}.
\begin{enumerate}
\item\label{sect4.1step1} We first establish the duality~\eqref{dualEFC} in the simpler case where $\mu$ has a second moment.
\item\label{sect4.1step2} We then construct a sequence of processes $(X_t^{(m)},t\geq 0)$ solution to~\eqref{SDEselec} with $f$ replaced by a certain smooth generating functions $f_m$, that are approximating $f$. We then show that $(X_t^{(m)},t\geq 0,m\geq 1)$ converges pointwise almost surely towards a certain process $(X_t^{(\infty)},t\geq 0)$ which satisfies the targeted duality relationship~\eqref{dualEFC}. See Lemma~\ref{approxXa}.
\item\label{sect4.1step3} We finally identify the process $(X_t^{(\infty)},t\geq 0)$ with the process $(X_t^{\min},t\geq 0)$ which is absorbed at the boundaries after reaching it. See Lemma~\ref{identificationXinfinity}.
\end{enumerate}
Our starting point is the following duality lemma which holds for $\Lambda$-WF processes with selection whose function $f$ is smooth.


\begin{lemm}\label{dualfinitemoment}
Assume $\mu(\infty)=0$ and $f''(1)=\sum_{n=2}^{\infty}n(n-1)\mu(n)<\infty$. Denote by $(X_t(x),t\geq 0)$ the solution to~\eqref{SDEselec}. Then for any $n\in \bbN$, and $x\in [0,1]$,
\begin{equation}\label{dualmoment1}
\bbE\left[X_t(x)^n\right]=\bbE\left[x^{N^{(n)}_t}\right].
\end{equation}
\end{lemm}

\begin{rema}\label{remKing}
This result has been observed by González et al. \cite[Lemma~2]{Gonzalesetal} in their study of branching processes with interaction. The moment duality~\eqref{dualmoment1} also holds true when $\Lambda$ gives mass at $0$. We check here that conditions for Ethier--Kurtz's theorem to hold are satisfied.
\end{rema}

\begin{proof}
Recall $\clA$ in~\eqref{generator} and $\clL^{c}$ in~\eqref{coalpart} the generators respectively of the $\Lambda$-WF process with no selection and the block counting process of the $\Lambda$-coalescent. For any $x\in [0,1]$ and $n\in \bbN$, set $g(x,n)= g_x(n)=g_n(x):=x^{n}$. The following identity is well-known, see e.g.~\cite[Lemma~2]{Gonzalesetal},
\begin{align}
\clA g_n(x)&=\clL^{c}g_x(n).
\end{align}
Recall $\clA^{\rms}$ in~\eqref{generatorWFs}. By assumption since $\mu(\infty)=0$, one has $\text{for all } x\in [0,1]$, $\clA^{\rms}g(x)=\clA g(x)+\mu(\bbN)x(f(x)-1)g'_n(x)$ and $\text{for all } x\in [0,1] \text{ and all } n\in \bbN$
\[
\mu(\bbN)x(f(x)-1)g'_n(x)=n\sum_{k=1}^{\infty}\big(g_x(n+k)-g_x(n)\big)\mu(k).
\]
This provides the duality at the level of the generators
\begin{multline}\label{expressionh}
h(x,n):=\clA^{\rms}g_n(x)=\clL g_x(n)\\
= \sum_{k=2}^{n}\left(x^{n-k+1}-x^{n}\right)\binom{n}{k}\lambda_{n,k}+\sum_{k=1}^{\infty}\left(x^{n+k}-x^{n}\right)n\mu(k).
\end{multline}

We now establish that the duality holds at the level of the semigroups. The process $(N_t^{(n)},t\geq 0)$ stays below a pure branching process $(Z_t,t\geq 0)$ with offspring measure $\mu$. By assumption $\mu$ admits a second moment, this entails in particular that for any time $T>0$, $\bbE(Z_T^{2})<\infty$, see e.g.~\cite[Chapter~3, Corollary~1 p.~111]{MR2047480}.\linebreak In particular, $(Z_t,t\geq 0)$ does not explode which ensures that $(N^{(n)}_t,t\geq 0)$ does not explode either. By Dynkin's formula for continuous-time Markov chains, since $g_x$ is bounded and the process does not explode, we see that
\[
\left(g\left(x,N_{t}^{(n)}\right)-\int_{0}^{t}h\left(x,N_s^{(n)}\right)\ddr s, t\geq 0\right)
\]
is a martingale. Since $\mu$ admits in particular a finite first moment, then the drift term $x\mapsto x(f(x)-1)$ is Lipschitz on $[0,1]$ and as noticed in Remark~\ref{Lipschitzrem}, it ensures that there is only one solution to the equation~\eqref{SDEselec} and that $1$ is absorbing for $X$ whenever it is reached. By applying Itô's formula to the process $(X_t(x),t\geq 0)$, we get that
\[
\left(g(X_{t}(x),n)-\int_{0}^{t}h(X_s(x),n)\ddr s, t\geq 0\right)
\]
is a local martingale. Since $g_n$ is bounded and $s\mapsto h(X_s(x),n)$ is bounded over finite time interval, the latter is a true martingale.

We now apply results of Ethier and Kurtz~\cite[Theorem~4.4.11, p.~192]{EthierKurtz}. Assume that $(X_t(x),t\geq 0)$ and $(N_t^{(n)},t\geq 0)$ are independent. Provided that the integrability assumption (4.50) of~\cite[Theorem~4.4.11]{EthierKurtz} is verified, Ethier and Kurtz's theorem (with in their notation $\alpha=\beta=0$) states that for all 
$
x\in [0,1], n\in \bbN, \bbE[x^{N_{t}^{(n)}}]\linebreak=\bbE[X_{t}(x)^{n}].
$
We check now assumption (4.50). Let $T>0$, clearly $\sup_{s,t\,\leq\,T}|g(X_{s},N_{t})|\linebreak\leq 1$ and it remains to see that the random variable $\sup_{s,t\,\leq\,T}|h(X_{s},N_{t})|$ is integrable. Recall the expression of $h(x,n)$ in~\eqref{expressionh}, one has for any $x\in[0,1]$ and $n\in \bbN$
\begin{align*}\label{integralform}
|h(x,n)|&\leq \sum_{k=2}^{n}\binom{n}{k}\lambda_{n,k}+n\mu(\bbN). 
\end{align*}
Recall the form of the $\lambda_{n,k}$'s in~\eqref{lambdank}. Simple binomial calculations provide that
\[
\sum_{k=2}^{n}\binom{n}{k}\lambda_{n,k}=\int_{0}^{1}\left(1-(1-z)^n-nz(1-z)^{n-1}\right)z^{-2}\Lambda(\ddr z).
\]
Setting $h(z)=1-(1-z)^n-nz(1-z)^{n-1}$, one checks $h'(u)=n(n-1)u(1-u)^{n-2}$ for all $u\in [0,1]$ and thus
\begin{align*}
\sum_{k=2}^{n}\binom{n}{k}\lambda_{n,k}&=\int_{0}^{1}z^{-2}\Lambda(\ddr z)\int_{0}^{z}n(n-1)u(1-u)^{n-2}\ddr u\\
&\leq n(n-1)\int_{0}^{1}z^{-2}\Lambda(\ddr z)\int_{0}^{z}u\ddr u\\
&=\frac{\Lambda([0,1])}{2}n(n-1).
\end{align*}
Therefore, for all $x\in [0,1]$ and all $n\in \bbN$, $|h(x,n)|\leq \frac{\Lambda([0,1])}{2}n(n-1)+\mu(\bbN)n$ and since $N_t\leq Z_t\leq Z_T$ for any $t\leq T$, one has almost surely
\[
\sup_{s,t\,\leq\,T}|h(X_{s},N_{t})|\leq \frac{\Lambda([0,1])}{2}Z_T^2+\mu(\bbN)Z_T:=\Gamma_T.
\]
The random variable $\Gamma_T$ is integrable since $\mu$ admits a second moment.
\end{proof}

We now go to step~\eqref{sect4.1step2}. Let $\mu$ be a finite measure on $\bar{\bbN}$ and $f$ be its generating function (possibly defective). For any $m\in \bbN$, recall $\bar{\mu}(m)=\mu(\{m,\,\ldots,\, \infty\})$ and $\mu_m$ defined in Lemma~\ref{Nnm}. For any $m\in \bbN$, set for any $x\in [0,1]$,
\begin{equation}\label{fm}
f_m(x):=\frac{1}{\mu_m(\bbN)}\sum_{k=1}^{\infty}x^{k}\mu_{m}(k)=\frac{1}{\mu(\bar{\bbN})}\left(\sum_{k=1}^{m-1}x^{k}\mu(k)+x^{m}\bar{\mu}(m)\right).
\end{equation}

\begin{lemm}\label{approxXa}
Let $(X_t^{(m)}(x),t\geq 0)$ be the unique strong solution to the SDE
\begin{multline}\label{sdem}
X^{(m)}_t(x)=x+\int_{0}^{t}\int_{0}^{1}\int_{0}^{1}z\left(\bbmun_{\left\{v\,\leq\,X^{(m)}_{s-}(x)\right\}}-X^{(m)}_{s-}(x)\right)\bar{\clM}(\ddr s,\ddr v, \ddr z) \\
-\mu_m(\bbN)\int_{0}^{t}X^{(m)}_{s-}(x)\big(1-f_m\left(X^{(m)}_{s-}(x)\right)\big)\ddr s.
\end{multline}
Then, for any $m\geq 1$,
\begin{equation}\label{order}
X_t^{(m+1)}\leq X_t^{(m)} \text{ for all } t\geq 0 \text{ almost surely}
\end{equation}
and the limiting process $(X^{(\infty)}_t,t\geq 0)$ defined by 
\[
X^{(\infty)}_t:=\underset{m\,\rightarrow\,\infty}{\lim}\!\! \downarrow X_t^{(m)} \text{ for all }t\geq 0,
\]
satisfies the duality relationship: for all $x\in [0,1)$, $n\in \bar{\bbN}$,
\begin{equation}\label{dualinfiny}
\bbE\left[X^{(\infty)}_t(x)^{n}\right]=\bbE\left[x^{N^{(n)}_t}\right].
\end{equation}
\end{lemm}

\begin{proof}
Note that for any $m\geq 1$, $f'_m(1-)<\infty$, therefore the SDE~\eqref{sdem} admits a unique strong solution. Moreover, $f_m(x)-f_{m+1}(x)=\bar{\mu}(m+1)(x^{m}-x^{m+1})\geq 0$ for any $x\in [0,1]$. By Theorem~\ref{comparisontheorem}, we see that~\eqref{order} holds true and thus $(X^{(\infty)}_t,t\geq 0) $ is well defined.

Recall Lemma~\ref{Nnm} and consider a monotone sequence of processes $(N^{(n)}_{m}(t),t\geq 0)$ with generator $\clL^m$ defined in~\eqref{genNm}, with splitting measure $\mu_m$ such that $\mu_{m}(k)=\mu(k)$ if $k\leq m-1$ and $\mu_m(m)=\bar{\mu}(m)$. Since $\mu_m$ admits a second moment, the duality relationship~\eqref{dualmoment1} holds and we have
\begin{equation}\label{dualm}
\bbE\left[x^{N_m^{(n)}(t)}\right]=\bbE\left[\left(X_t^{(m)}(x)\right)^{n}\right].
\end{equation}
By Lemma~\ref{Nnm}, $N_m^{(n)}(t)$ converges almost surely towards $N^{(n)}_t$ as $m$ goes to $\infty$. Hence, the identity~\eqref{dualinfiny} follows readily by taking limit on $m$.
\end{proof}

\begin{lemm}\label{identificationXinfinity}
The limit process $(X^{(\infty)}_t,t\geq 0)$ has the same law as $(X^{\min}_t,t\geq 0)$ the minimal solution of~\eqref{SDEselec}.
\end{lemm}

\begin{proof}
We first clarify the behavior of the process $(X^{(\infty)}_t(x),t\geq 0)$ when it reaches one of its boundaries $1$ or $0$. For the boundary $1$, since $\mu_m$ has a finite mean, \cite[Corollary~1.4 and Remark~1.5]{cdiEFC} apply and entail that the non-explosive processes $(N^{(\infty)}_m(t),t\geq 0)$ comes down from infinity. As observed in Section~\ref{backgroundWF}, this is equivalent to the fact that the dual process $(X_t^{(m)},t\geq 0)$ is getting absorbed at $1$ in finite time with positive probability. Note that by~\eqref{order}, $\tau_1^{(m)}:=\inf\{t\geq 0: X_t^{(m)}=1\}$ verifies $\tau_1^{(m+1)}\geq \tau_1^{(m)}$ and therefore
\[
\tau_1^{(\infty)}:=\inf\left\{t\geq 0: X^{(\infty)}_t=1\right\}\geq \underset{m\,\rightarrow\,\infty}{\lim}\uparrow \tau_1^{(m)}.
\]
Hence, since $1$ is absorbing for $(X_t^{(m)},t\geq 0)$ and $\tau_1^{(\infty)}\geq \tau_1^{(m)}$ a.s. If $\tau_1^{(\infty)}<\infty$ then 
\[
X^{(\infty)}_{t+\tau_1^{(\infty)}}=\underset{m\,\rightarrow\,\infty}{\lim} X^{(m)}_{t+\tau_1^{(\infty)}}=1 \text{ for any } t\geq 0
\]
a.s. On the event $\{\tau_1^{\infty}<\infty\}$, $(X^{(\infty)}_t,t\geq 0)$ is absorbed at its boundary $1$ in finite time, hence for any $t\geq 0$, $\bbP(X^{(\infty)}_t=1)=\bbP(\tau_1^{(\infty)}\leq t)$ and
\[
\bbP\left(\tau_1^{(m)}\leq t\right)=\bbP\left(X^{(m)}_t=1\right)\underset{m\,\rightarrow\,\infty}{\longrightarrow} \bbP(X^{(\infty)}_t=1)=\bbP(\tau_1^{(\infty)}\leq t),
\]
and thus $\tau_1^{(\infty)}=\underset{m\,\rightarrow\,\infty}{\lim}\uparrow \tau_1^{(m)}$ a.s.

For the boundary $0$, by taking $x=0$ in the duality relationship~\eqref{dualinfiny}, we see that $X_t^{(\infty)}(0)=0$ a.s. Hence, $0$ is necessarily absorbing.

We establish now that the process $(X^{(\infty)}_t,t>0)$ stopped at its first hitting time of the boundaries has the same law as the minimal process. This follows from uniform convergence of the generators. Recall $\clA^{\rms}$ the generator of the minimal solution $(X^{\min}_t,t\geq 0)$ to~\eqref{SDEselec}. For any $g\in C^2_c([0,1])$, $\clA^{\rms}g(x)=\clA g(x)+\mu(\bar{\bbN})x(f(x)-1)g'(x)$. Let $\clA^{\rms,(m)}$ be the generator of $(X_t^{(m)},t\geq 0)$. Since the jump parts of $\clA^{\rms}$ and $\clA^{\rms,(m)}$ are the same, and $\mu(\bar{\bbN})=\mu_m(\bar{\bbN})$ for any $m$, we have that
\begin{align}\label{diffgen}
\left\|\clA^{\rms,(m)}g-\clA^{\rms}g\right\|_{\infty}&=\mu(\bar{\bbN})\underset{x\,\in\, (0,1)}{\sup}\left|(f_m(x)-f(x))xg'(x)\right|.
\end{align}
For any $x\in [0,1]$,
\begin{align*}
0\leq f_m(x)-f(x)&=\frac{1}{\mu(\bar{\bbN})}\sum_{k\,\geq\,m}\left(x^{m}-x^{k}\right)\mu(k)\\
&\leq \frac{1}{\mu(\bar{\bbN})}x^{m}\sum_{j\,\geq\,0}(1-x^{j})\mu(j+m)\leq \frac{\bar{\mu}(m)}{\mu(\bar{\bbN})}x^m\leq x^m.
\end{align*}
Hence we see from~\eqref{diffgen} that $\|\clA^{\rms,(m)}g-\clA^{\rms}g\|_{\infty}\leq \sup_{x\in (0,1)}|x^{m+1}g'(x)|$. Since by assumption $g'$ has a compact support on $(0,1)$ and for any $x\in (0,1)$, $x^{m+1}\underset{m\,\rightarrow\,\infty}{\longrightarrow}0$, one has
\begin{equation}\label{convunifAm}
\|\clA^{\rms,(m)}g-\clA^{\rms}g\|_{\infty}\underset{m\,\rightarrow\,\infty}{\longrightarrow} 0.
\end{equation}
Moreover for large enough $m\geq 1$, $||\clA^{\rms,(m)}g\|_\infty\leq 1+\|\clA^{\rms}g\|_\infty$ and since $X_s^{(m)}\underset{k\,\rightarrow\,\infty}{\longrightarrow} X_s^{(\infty)}$ a.s. for any $s\geq 0$, $\clA^{\rms,(m)}g(X_s^{(m)})\underset{m\,\rightarrow\,\infty}{\longrightarrow} \clA^{\rms}g(X_s^{(\infty)})$ a.s. for any $s\geq 0$. Let $0\leq t_1\leq t_2\leq \ldots \leq t_n\leq s<t$ and $h_1,\,\ldots,\,h_n$ be some continuous functions defined on $[0,1]$. By Lemma~\ref{minimalprocess} applied to the process $(X_t^{(m)},t\geq 0)$, for any $g\in C_c^{2}([0,1])$, the process
\[
\left(g\left(X_{t\wedge \tau^{(m)}}^{(m)} \right)-\int_{0}^{t}\clA^{\rms,(m)}g\left(X_{s\,\wedge\,\tau^{(m)}}^{(m)}\right)\ddr s, t\geq 0\right)
\]
is a martingale. By applying Lebesgue's theorem, we have that
\begin{multline*}
\bbE_x\left[\left(g\left(X_{t\,\wedge\,\tau^{(\infty)}}^{(\infty)}\right)-g\left(X_{s\wedge \tau^{(\infty)}}^{(\infty)}\right)-\int_{s}^{t}\clA^{\rms}g\left(X_{r\wedge \tau^{(\infty)}}^{(\infty)}\right)\ddr r\right)\prod_{i=1}^{n}h_i\left(X_{t_i\,\wedge\,\tau^{(\infty)}}^{(\infty)}\right)\right]\\
\mkern-5mu{\begin{aligned}
&=\underset{m\,\rightarrow\,\infty}{\lim} \bbE_x\left[\left(g\left(X_{t\wedge \tau^{(m)}}^{(m)}\right)-g\left(X_{s\wedge \tau^{(m)}}^{(m)}\right)-\int_{s}^{t}\clA^{\rms,(m)}g\left(X_{r\wedge \tau^{(m)}}^{(m)}\right)\ddr r\right)\prod_{i=1}^{n}h_i\left(X_{t_i\wedge \tau^{(m)}}^{(m)}\right)\right]\\
&=0.
\end{aligned}}
\end{multline*}
This shows that the limiting process $(X^{(\infty)}_t,t\geq 0)$ stopped at time $\tau^{(\infty)}$ solves the martingale problem $(\MP)$. Lemma~\ref{minimalprocess} ensures that there is a unique solution to $(\MP)$, therefore $(X^{(\infty)}_t,t<\tau^{(\infty)})$ has the same law as the minimal process.
\end{proof}

Theorem~\ref{thmmomentdual1} is finally obtained by the combination of Lemma~\ref{approxXa} and Lemma~\ref{identificationXinfinity}. Indeed, the duality relationship~\eqref{dualEFC}
\[
\bbE\left[\left(X_t^{\min}(x)\right)^n\right]=\bbE\left[x^{N_t^{(n)}}\right] \text{ for any } x\in (0,1) \text{ and } n\in \bbN,
\]
is a direct consequence of~\eqref{dualinfiny}, By Lemma~\ref{monotonicity}, $N_t^{(n)}$ converges almost surely towards $N_t^{(\infty)}$, thus by dominated convergence theorem
\begin{equation}\label{entrancelawN0}
\bbP\left(X_t^{\min}(x)=1\right)=\underset{n\,\rightarrow\,\infty}{\lim} \bbE\left[\left(X_t^{\min}(x)\right)^n\right]=\underset{n\,\rightarrow\,\infty}{\lim} \bbE\left(x^{N_t^{(n)}}\right)=\bbE\left[x^{N_t^{(\infty)}}\right].
\end{equation}
Proof of Theorem~\ref{thmmomentdual1} is achieved.
\end{proof}


\subsection{Proofs of Theorems~\ref{Markovblockcounting} and~\ref{rec}}
These theorems state fundamental properties of the block counting process $(N_t,\linebreak t\geq 0):=(\#\Pi(t),t\geq 0)$ of any simple EFC process $(\Pi(t),t\geq 0)$. We highlight that apart in the proof of Corollary~\ref{inftyregularreflecting}, we will not make use of them in the next sections. We are going to show that $(N_t,t\geq 0)$ is a Feller Markov process with state space $\bar{\bbN}$ and that if it comes down from infinity then it is positive recurrent. The main tool will be the duality relationship~\eqref{dualEFC}.


\subsubsection{Proof of Theorem~\ref{Markovblockcounting}}

\begin{proof}
Let $(\Pi^{(n)}(t),t\geq 0)$ be the process following the coalescences and fragmentations involving only the first $n$ initial blocks of $(\Pi(t),t\geq 0)$, see Section~\ref{backgroundEFC} and~\cite[Lemma~3.3 and Lemma~3.4]{cdiEFC}. In particular, $\Pi^{(n)}(0):=(\Pi_1(0),\,\ldots,\,\Pi_n(0))$. When $n=\#\Pi(0)$, both processes $(N_t^{(n)},t\geq 0):=(\#\Pi^{(n)}(t),t\geq 0)$ and $(\#\Pi(t),t\geq 0)$ coincide. Without loss of generality, we enlarge the probability space on which $(\Pi(t),t\geq 0)$ is defined by assuming that the unique solution to~\eqref{SDEselec}, $(X_t^{\min}(x),t\geq 0)$, absorbed at the boundaries, is also defined on it and is independent of $(\Pi(t),t\geq 0)$. The Markov property of $(X_t^{\min}(x),t\geq 0)$ ensures that for any $s,t\geq 0$, conditionally given $X^{\min}_s(x)$, the random variable $X_{s+t}^{\min}(x)$ has the same law as $\tilde{X}_t^{\min}(X_{s}^{\min}(x))$, where the process $(\tilde{X}_t^{\min},t\geq 0)$ is an independent copy of $(X_t^{\min},t\geq 0)$ and is independent of $X_s^{\min}(x)$. Setting $n=\#\Pi(0)$ and applying Theorem~\ref{thmmomentdual1} and the duality relationship~\eqref{dualEFC}, we get for any $x\in [0,1]$ and any $s,t\geq 0$,
\begin{align*}
\bbE\left[x^{\#\Pi^{(n)}(s+t)}\right]&=\bbE\left[X_{s+t}^{\min}(x)^{n}\right]\\
&=\bbE\left[\bbE\left[\tilde{X}_{t}^{\min}\left(X_{s}^{\min}(x)\right)^{n}\middle|X_{s}^{\min}(x)\right]\right]\\
&=\bbE\left[X_s^{\min}(x)^{\#\Pi^{(n)}(t)}\right]\\
&=\bbE\left[\bbE\left[X_s^{\min}(x)^{\#\Pi^{(n)}(t)}\middle|\#\Pi^{(n)}(t)\right]\right]\\
&=\bbE\left[\bbE\left[x^{\#\tilde{\Pi}^{\left(\#\Pi^{(n)}(t)\right)}(s)}\middle|\#\Pi^{(n)}(t)\right]\right]\\
&=\bbE\left[x^{\#\tilde{\Pi}^{\left(\#\Pi^{(n)}(t)\right)}(s)}\right]
\end{align*}
where in the two last equalities, for any $m\in \bbN$, the random variable $\#\tilde{\Pi}^{(m)}(t)$ stands for an independent copy of $\#\Pi^{(m)}(t)$. We see finally that $\#\Pi(t+s)$ has the same distribution as $\#\tilde{\Pi}^{(\#\Pi(t))}(s)$. The process $(\#\Pi(t),t\geq 0)$ is therefore Markovian.

We now establish the Feller property. Recall $\bar{\bbN}$ the one-point compactification of $\bbN$ and $C(\bar{\bbN})$ the space of continuous functions defined on $\bar{\bbN}$. Rewriting the limit~\eqref{entrancelawN0} in terms of the processes $(\#\Pi^{(n)}(t),t\geq 0)$ and $(\#\Pi(t),t\geq 0)$, we see that if $\#\Pi(0)=\infty$ a.s. then
\begin{equation}\label{entrancelawN}
\underset{n\,\rightarrow\,\infty}{\lim} \bbE_n\left(x^{\#\Pi(t)}\right)=\underset{n\,\rightarrow\,\infty}{\lim} \bbE\left(x^{\#\Pi^{(n)}(t)}\right)=\bbP(X_t^{\min}(x)=1)=\bbE\left[x^{\#\Pi(t)}\right].
\end{equation}
Recall
\[
g_x(n)=x^{n}\quad \text{and set}\quad x^{\infty}:=\underset{n\,\rightarrow\,\infty}{\lim} x^{n}=\bbmun_{\{x=1\}}.
\]
We have just established in~\eqref{entrancelawN} that $n\mapsto \bbE_n(g_x(\#\Pi(t)))$ is continuous at $\infty$. The subalgebra $A$ of $C(\bar{\bbN})$ generated by the linear combinations of the maps $\{n\mapsto g_x(n), x\in [0,1]\}$, is separating $C(\bar{\bbN})$. Moreover, for any $n\in \bar{\bbN}$, there is a function $g$ in $A$ such that $g(n)\neq 0$. By the Stone--Weierstrass theorem, $A$ is dense in $C(\bar{\bbN})$ for the uniform norm, since $n\mapsto \bbE_n(g(\#\Pi(t))) $ is continuous for any $g\in A$, this holds true for any $g\in C(\bar{\bbN})$ and the semigroup of $(\#\Pi(t),t\geq 0)$ maps $C(\bar{\bbN})$ to $C(\bar{\bbN})$. It remains to verify the continuity of the semigroup at $0$. This is a direct application of the duality relationship~\eqref{momentdual1}, since $X^{\min}$ is right-continuous. Finally, the process $(\#\Pi(t),t\geq 0)$ is Feller.
\end{proof}

\subsubsection{Proof of Theorem~\ref{rec}}

\begin{proof}
 Our objective is to show that when the block counting process comes down from infinity, it is positive recurrent. We shall use the duality relationship~\eqref{dualEFC}. Denote by $(Y_t(x),t\geq 0)$ the neutral $\Lambda$-Wright--Fisher process, that is to say, the unique solution to~\eqref{SDEselecintro} with $f(x)=1$ for all $x\in [0,1]$. As recalled in the introduction, when the $\Lambda$-coalescent process comes down, the process $(Y_t(x),t\geq 0)$ has a positive probability to be absorbed at $0$. Moreover, for any $m\geq 1$, $x(f_m(x)-1)\leq 0$, and by applying the comparison theorem between $(X_t^{(m)}(x),t\geq 0)$ and $(Y_t(x),t\geq 0)$, we get $X_t^{(m)}(x)\leq Y_t(x)$ for all $t\geq 0$ almost surely. Hence, letting $m$ converge to $\infty$, we see that $X^{\min}_t(x)\leq Y_t(x)$ for all $t\geq 0$ almost surely, which ensures that the process $(X^{\min}_t(x),t\geq 0)$ hits $0$ with positive probability. Moreover, according to Theorem~\ref{thmmomentdual1}-i), since $(N_t^{(\infty)},t\geq 0)$ comes down from infinity, the process $(X^{\min}_t(x),t\geq 0)$ hits $1$ with positive probability. Finally, since process $(X_t^{\min},t\geq 0)$ is a positive supermartingale, it converges almost surely as $t$ goes to $\infty$ towards one of its absorbing boundaries $0$ or $1$. On the event $\{\tau_1<\tau_0\}$, $\underset{t\,\rightarrow\,\infty}{\lim} X_t^{\min}=1$ a.s and similarly on the event 
$
\{\tau_0<\tau_1\}, \underset{t\,\rightarrow \,\infty}{\lim} X_t^{\min}=0
$
a.s. Thus, by the duality relationship~\eqref{dualEFC},
\[ 
\underset{t\,\rightarrow\,\infty}{\lim} \bbE\left[x^{N_t^{(n)}}\right]=\bbP_x(\tau_1<\tau_0).\qedhere
\]
\end{proof}


\section{Proof of Theorem~\ref{thmmomentduality2}: construction of the extended process \texorpdfstring{$(X^{\rmr}_t,t\geq 0)$}{(Xrt,t>= 0)}}\label{sec2}
The assumption $\Lambda(\{0\})=0$ will play an important role in this section. We consider from now on such a coalescence measure $\Lambda$ (also with no mass at $1$) and a splitting measure $\mu$ without mass at $\infty$. The only process solution to~\eqref{SDEselec} whose existence is clear is the minimal process. It will be therefore necessary to construct the process $X^{\rmr}$ whose boundary $1$ is not necessarily absorbing. We stress that our study focuses on extensions of the minimal process at the barrier $1$. The boundary $0$ will be always absorbing. We will construct $(X_t^{\rmr}, t\geq 0)$ as limit of certain processes $(X_t^{\lambda},t\geq 0,\lambda\linebreak>0)$ whose boundary $1$ is an entrance. Recall that we wish to establish the second duality relationship~\eqref{dualityEFCref}: for any $x\in (0,1)$ and $n\in \bbN$,
\[
\bbE\left[X_t^{\rmr}(x)^n\right]=\bbE\left[x^{N_t^{\min,(n)}}\right].
\]
Once this identity obtained, the correspondences given in Table~\ref{correspondance} and~\ref{correspondancereg} will follow. The proof of Theorem~\ref{thmmomentduality2} is deferred. We explain first the strategy and establish several lemmas.
\begin{itemize}
\item[-] \textbf{\emph{Step 1:}} Let $\mu$ be a splitting measure with no mass at $\infty$ and $f$ its associated probability generating function. We consider a family of $\Lambda$-Wright--Fisher processes with selection, $(X^{\lambda}_t,t\geq 0)$ indexed by $\lambda>0$. Each has its selection mechanism driven by the defective function $f^{\lambda}$ associated to $\mu^{\lambda}=\mu+\lambda\delta_{\infty}$. We then establish a duality relationship between this process $(X^{\lambda}_t,t\geq 0)$ and the block counting process $(N_t^{\lambda},t\geq 0)$ of a simple EFC process whose splitting measure is $\mu^{\lambda}$. We show that under the assumption $\Lambda(\{0\})=0$, $(X^{\lambda}_t,t\geq 0)$ has boundary $1$ \emph{entrance}. This is the aim of Lemma~\ref{recextenslambda}.
\item[-] \textbf{\emph{Step 2:}} We study the dual processes $(N_t^{\lambda},t\geq 0)$ as $\lambda$ goes to $0$ and establish that they converge as $\lambda$ goes to $0$ towards $(N_t^{\min},t\geq 0)$, the block counting process with splitting measure $\mu$ and coalescence measure $\Lambda$ that is stopped after it has reached the boundary $\infty$. This is the aim of Lemma~\ref{Nlambda}.
\item[-] \textbf{\emph{Step 3:}} The convergence of the processes $(N_t^{\lambda},t\geq 0)$ as $\lambda$ goes to $0$ shown in Step 2 entails the convergence of processes $(X^{\lambda}_t,t\geq 0)$ considered in Step 1. We study the limit process called $(X^{\rmr}_t,t\geq 0)$, establish the duality relationship~\eqref{dualityEFCref} and verify that this is an extension of the minimal $\Lambda$-WF process with selection driven by $f$. This is the aim of Lemma~\ref{recextensX}.
\item[-] \textbf{\emph{Step 4:}} We study the possible behaviors at the boundary $1$ of the extended process $(X_t^{\rmr},t\geq 0)$ from the duality~\eqref{dualityEFCref}. We first establish from the duality a correspondence between boundaries that are non-absorbing and those accessible in Lemma~\ref{correspond}. The correspondences~\eqref{theo2.7.i} to~\eqref{theo2.7.iv} stated in Theorem~\ref{thmmomentduality2} are then established. Recall that when the process started from $1$ is degenerate\linebreak at $1$, and further $1$ is inaccessible, we say that the boundary $1$ is natural. In the case that $1$ is accessible, the boundary is an exit. When the process started from $1$ leaves $1$ and never returns to it again almost surely, $1$ is an entrance. Finally, when the process started from $1$, leaves it almost surely and returns to it with positive probability (i.e. $1$ is accessible), the boundary is regular non-absorbing.
\end{itemize}

%\begin{proof}[{Proof of Theorem~\ref{thmmomentduality2}}]
\begin{step}%\textbf{Step 1}.
Let $\Lambda$ be a coalescence measure with no atom at $0$ nor $1$. Let $\mu$ be a finite measure on $\bbN$ and denote by $f$ its probability generating function. Fix $\lambda>0$. We denote by $(N_t^{\lambda,(n)},t\geq 0)$ the block counting process started from $n\in \bar{\bbN}$ with coalescence measure $\Lambda$ and splitting measure $\mu^{\lambda}$ defined such that $\mu^{\lambda}(k)=\mu(k)$ for any $k\in \bbN$ and $\mu^{\lambda}(\infty)=\lambda$. Let $f^{\lambda}$ be the defective probability generating function associated to $\mu^{\lambda}$. For any $x\in [0,1]$, 

\[
f^{\lambda}(x)=\sum_{k=1}^{\infty}x^{k}\frac{\mu(k)}{\lambda+\mu(\bbN)}=\frac{\mu(\bbN)}{\lambda+\mu(\bbN)}f(x)\quad\text{and}\quad f^{\lambda}(1)-1=-\frac{\lambda}{\mu(\bbN)+\lambda}<0.
\]

Recall the SDE~\eqref{SDEselec}. Let $(X^{\lambda}_t,t\geq 0)$ be the minimal process solution to the equation
\begin{multline}\label{sdelambda}
X^{\lambda}_t(x)=x+\int_{0}^{t}\int_{0}^{1}\int_{0}^{1}z\left(\bbmun_{\left\{v\leq X^{\lambda}_{s-}(x)\right\}}-X^{\lambda}_{s-}(x)\right)\bar{\clM}(\ddr s,\ddr v, \ddr z)\\
-\mu^{\lambda}(\bar{\bbN})\int_{0}^{t}X^{\lambda}_{s}(x)\left(1-f^{\lambda}\left(X^{\lambda}_{s}(x)\right)\right)\ddr s.
\end{multline}

\begin{lemm}\label{recextenslambda}
The process $(X_t^{\lambda}(x),t\geq 0)$ verifies the duality relationship
\begin{equation}\label{duallambda}
\bbE\left[X^{\lambda}_t(x)^n\right]=\bbE\left[x^{N_t^{\lambda,(n)}}\right],
\end{equation}
for any $n\in \bar{\bbN}$, any $x\in [0,1)$ and $t\geq 0$. Moreover, $(X_t^{\lambda},t\geq 0)$ has $1$ as an entrance boundary, and the entrance law at $1$ is characterized via its moments by
\begin{multline}\label{entrancelawlambda}
\bbE\left[X^{\lambda}_t(1)^n\right]:=\underset{x\,\rightarrow\,1-}{\lim} \bbE\left[x^{N_t^{\lambda,(n)}}\right]\\
=\bbP\left(N_t^{\lambda,(n)}<\infty\right),\text{ for any } t\geq 0 \text{ and any } n\in \bbN.
\end{multline}
The semigroup of $(X_t^{\lambda}(x),t\geq 0, x\in [0,1])$ satisfying~\eqref{duallambda} and~\eqref{entrancelawlambda} is Feller.
\end{lemm}

\begin{rema}
When the measure $\mu$ is a Dirac mass at $\infty$, $\mu=\lambda\delta_{\infty}$, the drift term in the stochastic equation~\eqref{sdelambda} reduces to $-\lambda\int_{0}^{t}X_s^{\lambda}(x)\ddr s$. The process $(X_t^{\lambda}(x),t\geq 0)$ corresponds to a $\Lambda$-Wright--Fisher process with no selection but unilateral mutation from type $a$ to $A$ at rate $\lambda$. A side consequence of Lemma~\ref{recextenslambda} is that the boundary $1$ is an entrance (thus inaccessible) for the $\Lambda$-Wright--Fisher process with unilateral mutation when $\Lambda(\{0\})=\Lambda(\{1\})=0$.
\end{rema}

\begin{proof}
Let $(X_t^{\lambda,\min},t\geq 0)$ be the process solution to~\eqref{sdelambda} that is stopped when reaching boundary $1$. Namely 
\[
\left(X_t^{\lambda,\min},t\geq 0\right):=\left(X_{t\,\wedge\,\tau_1}^{\lambda},t\geq 0\right).
\]
By Theorem~\ref{thmmomentdual1}, the process satisfies the identity $\bbE[X_t^{\lambda,\min }(x)^n]=\bbE[x^{N_t^{\lambda,(n)}}]$ for any $x\in [0,1)$ and $n\in \bbN$. Recall that the key Lemma~\ref{key} ensures that for any $\lambda>0$, the process $(N^{\lambda}_t,t\geq 0)$ with splitting measure $\mu^{\lambda}$ has $\infty$ as an exit boundary (i.e. $\infty$ is accessible absorbing). Moreover, by Lemma~\ref{monotonicity}, $N_t^{\lambda,(\infty)}=\underset{n\,\rightarrow\,\infty}{\lim} N_t^{(n)}$ a.s. Therefore 
\[
\bbE\left[x^{N_t^{\lambda,(\infty)}}\right]=\bbP\left(X_t^{\lambda,\min }(x)=1\right)=0
\]
and $1$ is inaccessible for
\[
\left(X_t^{\lambda,\min },t\geq 0\right)\quad \text{and for}\quad\left(X_t^{\lambda},t\geq 0\right).
\]
Thus, one has $\bbE[X^{\lambda}_t(x)^n]=\bbE[x^{N_t^{\lambda,(n)}}]$ and letting $x$ go towards $1$, we see that
\[
\underset{x\,\rightarrow\,1-}{\lim}  \bbE\left[X^{\lambda}_t(x)^n\right]=\bbP\left(N_t^{\lambda,(n)}<\infty\right)\in (0,1).
\]
This characterizes an entrance law at boundary $1$ for the process $(X_t^{\lambda},t\geq 0)$. We establish the Feller property of the extended semigroup of $(X^{\lambda}_t,t\geq 0)$ on $[0,1]$. Recall $g_n(x)=g_x(n)=x^{n}$. Plainly by the duality relationship~\eqref{duallambda}, if one denotes by $(P^{\lambda}_t)$ the semigroup of $(X_t^{\lambda},t\geq 0)$, we see that
\[
P^{\lambda}_tg_n:x\mapsto \bbE\left[X_t^{\lambda}(x)^n\right]=\bbE\left[x^{N_t^{\lambda,(n)}}\right]
\]
is continuous on $[0,1)$. Thus, for any polynomial function $h$ on $[0,1]$, $x\mapsto \bbE[h(X_t^{\lambda}(x))]$ is continuous on $[0,1]$. By the Weierstrass theorem, if $g\in C([0,1])$, one can find a sequence of polynomial functions $(h_n)$ such that $h_n\underset{n\,\rightarrow\, \infty}{\longrightarrow} g $ uniformly. The following routine calculation establishes the continuity of $P^{\lambda}_t g$: for any $x,y\in [0,1]$,
\begin{multline*}
\left\lvert P^{\lambda}_tg(x)-P_t^{\lambda}g(y)\right\lvert\\
\begin{aligned}
&=\left\lvert P_t^{\lambda}g(x)-P_t h_n(x)+P_t^{\lambda}h_n(x)-P_t^{\lambda}h_n(y)+P_t^{\lambda}h_n(y)-P_t^{\lambda}g(y)\right\lvert\\
&\leq 2\|g-h_n\|_\infty+\left|P_t^{\lambda}h_n(x)-P_t^{\lambda}h_n(y)\right|.
\end{aligned}
\end{multline*}
Since $P_t^{\lambda} h_n$ is continuous on $[0,1]$, if one let $x$ tend to $y\in [0,1]$ and then $n$ to $\infty$, we get
\[
\underset{x\,\rightarrow\,y}{\limsup}\left\lvert P_t^{\lambda}g(x)-P_t^{\lambda}g(y)\right\lvert\leq \|g-h_n\|_\infty\underset{n\,\rightarrow\,\infty}{\longrightarrow}0,
\]
which allows us to conclude that $P_t^{\lambda}$ maps $C([0,1])$ into $C([0,1])$. We now check the strong continuity at $0$ of the semigroup $P_t^{\lambda}$. Since it is Feller, it is sufficient to check the pointwise continuity, and by the Weierstrass theorem, we can focus on the functions $g_n:x\mapsto x^{n}$, namely we need to show $\bbE[X^{r}_t(x)^{n}]\underset{t\,\rightarrow\,0+}{\longrightarrow} x^{n}$. The latter follows readily from the duality~\eqref{duallambda} and the right-continuity of $(N_t^{\lambda},t\geq 0)$.
\end{proof}\let\qed\relax
\end{step}

\vspace*{4pt}
\begin{step}
We now study the dual process $(N_t^{\lambda}, t\geq 0) $. Recall that by assumption $\Lambda(\{0\})=0$. The key Lemma~\ref{key} plays again a central role in the proof of the following lemma. Recall that $(N^{\min,(n)}_t,t\geq 0)$ denotes the block counting process started from $n\in \bbN$ and absorbed at $\infty$ whenever it reaches it.

\begin{lemm}\label{Nlambda}
Let $(N_t^{(n)},t\geq 0)$ be a block counting process with coalescence measure $\Lambda$ (with no mass at $0$) and splitting measure $\mu$ (with no mass at $\infty)$. There exists on the same probability space, block counting processes $(N_t^{\lambda,(n)},t\geq 0,\lambda>0)$ started from $n\in \bbN$ with splitting measure $\mu^{\lambda}:=\mu+\lambda\delta_\infty$ and the same coalescence measure $\Lambda$, such that if $\lambda'\leq \lambda$ then
\begin{equation}\label{orderlambda}
N_t^{\lambda',(n)}\leq N_t^{\lambda,(n)} \text{ for all }t\geq 0 \text{ a.s.}
\end{equation}
Almost surely for any $n\in \bbN$, $\underset{\lambda\,\rightarrow\,0+}\lim \!\!\! \downarrow N_t^{\lambda,(n)}=N_t^{\min,(n)} \text{ for all } t\geq 0$ and
\begin{multline*}
\zeta_{\infty}^{\lambda}\underset{\lambda\, \rightarrow\,0+}{\longrightarrow} \zeta_{\infty}\text{ a.s.  where }\zeta_{\infty}^{\lambda}:=\inf
\left\{t>0: N_t^{\lambda,(n)}=\infty\right\}\\
\text{and}\quad \zeta_{\infty}:=\inf\left\{t>0: N_{t-}^{(n)}=\infty\right\}.
\end{multline*}
\end{lemm}

\begin{proof}
We work at the level of the partition-valued processes. We are going to construct a collection of EFC processes $(\Pi^{\lambda}(t),t\geq 0,\lambda>0)$ with coalescence measure $\Lambda$ and splitting measures $\mu^{\lambda}=\mu+\lambda \delta_{\infty}$. Let $\mu_{\Frag}$ be an exchangeable measure on $\clP_{\infty}$ such that $\mu_{\Frag}(\pi;\#\pi^{f}-1=k)=\mu(k)$ for all $k\in \bbN$. Recall the dynamics of the process $(\#\Pi(t),t\geq 0)$ explained in Section~\ref{backgroundEFC}. Let $\PPP_F$ and $\PPP_C$ be the Poisson point processes governing respectively fragmentations and coalescences. To incorporate a mass $\lambda$ at $\infty$ in the splitting measures, we consider an additional term of fragmentation along a partition with infinitely many blocks. In order to do so, let $\rmm:=(m_1,m_2,\,\ldots)$ be a positive sequence such that $\sum_{i=1}^{\infty}m_i=1$ and $m_i>m_{i+1}>0$ for all $i\geq 1$. Denote by $\rho_{\rmm}$ the law of the paint-box partition whose ranked asymptotic frequencies are given by $\rmm$ and note that $\#\pi=\infty$, for $\rho_{\rmm}$-almost every partition $\pi$. Consider now $\PPP^{1}:=\sum_{i\,\geq\,1}\delta_{(t^{1}_i,\pi_i,j)}$ an independent Poisson point process, with intensity $\ddr t\otimes \rho_{\rmm}\otimes \#$, where we recall that $\#$ is the counting measure. Let $(\PPP^{\lambda}, \lambda>0)$ be the images of $\PPP^{1}$ by the map $t\mapsto \lambda t$. They are Poisson point processes with intensity $\lambda \ddr t\otimes \rho_{\rmm}\otimes \#$, such that if $\lambda\geq \lambda'$, then $t_1^{\lambda}\leq t_1^{\lambda'}$ almost surely where $t_1^{\lambda}$ denotes the first atom of time of $\PPP^{\lambda}$. Let $(\Pi^{\lambda}(t),t\geq 0)$ for $\lambda>0$ be the simple EFC processes built from the Poisson point process $\PPP^{\lambda}_{F}=\PPP_F+\PPP^{\lambda}$ and $\PPP_C$. For any $\lambda>0$, the fragmentation measure is $\mu^{\lambda}_{\Frag}:=\mu_{\Frag}+\lambda \rho_{\rmm}$, and thus the splitting measure is $\mu^{\lambda}$ such that $\mu^{\lambda}(k)=\mu(k)$ for all $k\in \bbN$ and $\mu^{\lambda}(\infty)=\lambda \rho_{\rmm}(\pi; \#\pi=\infty)=\lambda$. Recall the process $(\Pi^{\lambda,(n)}(t),t\geq 0)$ defined in Section~\ref{backgroundEFC}. We set $(N_t^{\lambda,(n)},t\geq 0):=(\#\Pi^{\lambda,(n)}(t),t\geq 0)$ and $\zeta_{\infty}^{\lambda}:=\inf\{t>0: N_{t-}^{\lambda,(n)}\text{ or } N_{t}^{\lambda,(n)}=\infty\}$. By Lemma~\ref{key}, all processes $N^{\lambda}$ have boundary $\infty$ exit. We see by construction that almost surely for any $\lambda>\lambda'>0$ and any $t\geq 0$,
\begin{equation}\label{comparisonNlambda}
N^{\lambda,(n)}_t\geq N_t^{\lambda',(n)}\geq N_t^{(n)}
\text{ and } \zeta_{\infty}^{\lambda}\leq \zeta_{\infty}^{\lambda'}\leq \zeta_{\infty}.
\end{equation}
For any $t\geq 0$, set 
\[
N_t^{0,(n)}:=\underset{\lambda\,\rightarrow\,0+}{\lim}\downarrow N_t^{\lambda}\quad\text{a.s and}\quad\zeta^0_{\infty}:=\inf\left\{t>0: N_{t-}^{0,(n)}=\infty\right\}.
\]
By~\eqref{comparisonNlambda}, almost surely for any $\lambda>0$, $\zeta^{0}_{\infty}\geq \zeta_{\infty}^{\lambda}$ a.s. Hence, on the event $\{\zeta_{\infty}^{0}<\infty\}$, for any $t>0$, $t+\zeta_{\infty}^{0}\geq t+\zeta_{\infty}^{\lambda}$ and Lemma~\ref{key} yields that $N^{\lambda,(n)}_{t+\zeta_{\infty}^{0}}=\infty$ for all $t\geq 0$ a.s. Therefore
\[
N^{0,(n)}_{t+\zeta_{\infty}^{0}}=\underset{\lambda \,\rightarrow\,0+}{\lim} N^{\lambda,(n)}_{t+\zeta_{\infty}^{0}}=\infty \text{ a.s.}
\]
The process $(N_t^{0,(n)},t\geq 0)$ is thus absorbed at $\infty$ after its first explosion time $\zeta_\infty^0$. By construction, $(N_t^{0,(n)},t<\zeta^{0}_{\infty})$ has the same dynamics as the block counting process $(N_t^{(n)},t<\zeta_{\infty})$ whose splitting measure is $\mu$ and coalescence measure is $\Lambda$. By the uniqueness of the minimal continuous-time Markov chain with generator $\clL$, $(N_t^{0,(n)},t\geq 0)$ and the stopped process $(N_t^{\min,(n)},t\geq 0):=(N_{t\wedge \zeta_{\infty}}^{(n)},t\geq 0)$ have the same law. On the other hand, \eqref{comparisonNlambda} entails that almost surely for any $t\geq 0$,
\begin{equation}\label{comparisonNN0}
N_t^{0,(n)}\geq N_t^{(n)} \text{ and } \zeta_{\infty}^0\leq \zeta_{\infty}.
\end{equation}
\pagebreak

\noindent Since $\zeta_{\infty}^0$ and $\zeta_{\infty}$ have the same law, they coincide almost-surely. A similar reasoning using~\eqref{comparisonNN0} entails that $(N_t^{0,(n)},t\geq 0)$ and $(N^{(n)}_{t\,\wedge\,\zeta_\infty},t\geq 0)$ are equal almost-surely.
\end{proof}\let\qed\relax
\end{step}
\vspace*{4pt}
%\textbf{Step 3}:

\begin{step}
 We study now some extensions of the minimal process, solution to the SDE~\eqref{SDEselec}, whose drift term satisfies $f(1)=1$. We define our extension $X^{\rmr}$ by looking at the limit arising in the processes $(X^{\lambda},\lambda>0)$ when the mutation rate $\lambda$ gets very low. The duality relationship~\eqref{dualityEFCref} will be shown at the same time as the convergence. The convergence of the sequence of processes with mutation $(X^{\lambda},\lambda>0)$ towards $X^{\rmr}$ holds in Skorokhod sense. We will see that the convergence holds also pointwise almost surely in Lemma~\ref{almostsureinx}.


\begin{lemm}[Extension of $(X^{\min}_t,t\geq 0)$ after reaching $1$]\label{recextensX}
The Markov processes $(X^{\lambda}_t(x),t\geq 0, x\in [0,1])$ converge as $\lambda$ goes to $0$, in the Skorokhod sense towards a Feller process $(X_t^{\rmr}(x),t\geq 0,x\in [0,1])$ valued in $[0,1]$, which extends the minimal solution of~\eqref{SDEselec}, and whose semigroup satisfies : for any $n\in \bbN$
\begin{align}
\bbE\left[X^{\rmr}_t(x)^n\right] &=\bbE\left[x^{N^{\min,(n)}_t}\right]
\text{ for any } x\in [0,1),\label{dualabso}
\\
\intertext{and}
\bbE\left[X^{\rmr}_t(1)^n\right]\, &{:=}\underset{x\,\rightarrow\,1-}{\lim}\bbE\left[x^{N^{\min,(n)}_t}\right]=\bbP_n(\zeta_\infty>t).\label{entrancelawXr}
\end{align}
\end{lemm}

\begin{rema}\label{remwhynotkingman}
In order to establish Lemma~\ref{recextensX}, we require the assumption that the measure $\Lambda$ gives no mass to $0$. We shall indeed use the fact that when there is no Kingman component, the processes $N^{\lambda}$ have all their boundary $\infty$ as exit. This is not the case when $\Lambda(\{0\})=c_{\rmk}>0$ for which processes $N^{\lambda}$ may have boundary $\infty$ regular, see Remark~\ref{remkey}.
\end{rema}

\begin{proof}
In order to ease the reading we outline the scheme of the proof. The strategy is similar to that in the proof of Lemma~\ref{recextenslambda}:
\begin{enumerate}
\item\label{proof1} We first show that for any fixed $t\geq 0$, the random variables $(X_t^{\lambda}(x),\linebreak x \in[0,1])$ converge in law as $\lambda$ goes to $0$ through convergence of their moments. The limiting random variables $(X_t^{\rmr}(x),x\in [0,1])$ satisfy the identities~\eqref{dualabso} and~\eqref{entrancelawXr}.
\item\label{proof2} We establish that the semigroup of $(X^{\lambda}_t(x),t\geq 0, x\in [0,1])$ converges as $\lambda$ goes to $0$ uniformly. We denote by $P_t^{\rmr}$ the limiting operator and verifies that it is a Feller semigroup. The associated process is denoted by $(X^{\rmr}_t(x),\linebreak t\geq 0,x\in [0,1])$.
\item\label{proof3} We then show that $(X^{\rmr}_t(x),t\geq 0,x\in [0,1])$ extends the minimal solution to~\eqref{SDEselec} after reaching $1$.
\end{enumerate}
\eqref{proof1} By letting $\lambda$ go towards $0$ in the duality relationship~\eqref{duallambda}, and recalling the almost sure convergence of $N_t^{\lambda,(n)}$ towards $N_t^{\min,(n)}$, see Lemma~\ref{Nlambda}, we see by the dominated convergence theorem, that
\[
\underset{\lambda\,\rightarrow\,0}{\lim}\ \bbE\left[X_t^{\lambda}(x)^{n}\right]=\bbE\left[x^{N_t^{\min,(n)}}\right].
\]
Recall that the convergence in law of random variables valued in $[0,1]$ is characterized by the convergence of the entire moments. Therefore, the $X_t^{\lambda}(x)$'s are converging in law as $\lambda$ goes to $0$ towards some random variable $X_t^{\rmr}(x)$ whose law is characterized by its entire moments defined by $\bbE[(X^{\rmr}_t(x))^{n}]:=\bbE[x^{N_t^{\min,(n)}}]$. In particular, $X_t^{\rmr}(x)$ satisfies the targeted duality relationship~\eqref{dualabso}. Similarly by letting $\lambda$ go to $0$ in the identity~\eqref{entrancelawlambda}, recalling Lemma~\ref{Nlambda} and the fact that $\zeta_\infty^{\lambda}$ goes to $\zeta_{\infty}$ a.s., we also get that the identity~\eqref{entrancelawXr} hold.

\noindent\eqref{proof2} Recall that by Lemma~\ref{recextenslambda}, $(X^{\lambda}_t(x),t\geq 0)$ is a Feller process and that we denote its semigroup by $(P_t^{\lambda},t\geq 0)$. Let $g_n(x)=x^{n}$ for any $x\in [0,1]$ and $n\in \bbN$. By the duality relationship~\eqref{duallambda}, $P_t^{\lambda}g_n(x)=\bbE[x^{N^{(n),\lambda}_{t}}]$. We check that
\begin{equation}\label{unifconvsemigroup}
\left\|P_t^{\lambda}g_n-P_t^{\rmr}g_n\right\|_\infty=\sup_{x\,\in\,[0,1]}\bbE\left[x^{N_t^{\min,(n)}}-x^{N^{(n),\lambda}_{t}}\right]\underset{\lambda\,\rightarrow\,0}{\longrightarrow} 0.
\end{equation}
Arguments are adapted from those in~\cite[Section~7]{zbMATH07055671}. For any $x\in [0,1]$,
\begin{multline*}
\bbE\left[x^{N_t^{\min,(n)}}-x^{N^{(n),\lambda}_t}\right]\\
\begin{aligned}
&=\bbE\left[\left(x^{N^{\min,(n)}_{t}}-x^{N_t^{\lambda,(n)}}\right)1_{\left\{N_t^{\min,(n)}\,\leq\,N_t^{\lambda,(n)}\,<\,\infty\right\}}\right]\\
&\qquad +\bbE\left[\left(x^{N^{\min,(n)}_{t}}-x^{N_t^{\lambda,(n)}}\right)1_{\left\{N_t^{\min,(n)}<\,\infty, N_t^{\lambda,(n)}=\infty\right\}}\right]\\
&\leq \bbE\left[\left(x^{N^{\min,(n)}_{t}}-x^{N_t^{\lambda,(n)}}\right)1_{\left\{N_t^{\min,(n)}\,\leq\,N_t^{\lambda,(n)}\,<\,\infty\right\}}\right]+2\bbP_n\left(\zeta_{\infty}>t\geq \zeta_{\infty}^{\lambda}\right).
\end{aligned}
\end{multline*}
Recall that by Lemma~\ref{Nlambda}, $\zeta_{\infty}^{\lambda}\underset{\lambda\, \rightarrow\,0+}{\longrightarrow} \zeta_{\infty}$ a.s, thus $\bbP_n(\zeta_{\infty}>t\geq \zeta_{\infty}^{\lambda})\underset{\lambda\,\rightarrow\,0}{\longrightarrow} 0$. It remains to study the uniform convergence on the event $\{N_t^{\min,(n)}\leq N_t^{\lambda,(n)}<\infty\}$. Plainly,
\begin{multline}\label{upperboundunif}
\underset{x\,\in\,[0,1]}{\sup} \bbE\left[\left(x^{N^{\min,(n)}_{t}}-x^{N_t^{\lambda,(n)}}\right)1_{\left\{N_t^{\min,(n)}\leq N_t^{\lambda,(n)}<\infty\right\}}\right]\\
\leq \bbE\left[\underset{x\,\in\,[0,1]}{\sup}\left(x^{N^{\min,(n)}_{t}}-x^{N_t^{\lambda,(n)}}\right)1_{\left\{N_t^{\min,(n)}\,\leq\,N_t^{\lambda,(n)}<\infty\right\}}\right].
\end{multline}
By Lemma~\ref{Nlambda}, $N_t^{\lambda,(n)} \underset{\lambda\,\rightarrow\,0}{\longrightarrow} N^{\min,(n)}_t$ a.s. On the event $\{N_t^{\min,(n)}\leq N_t^{\lambda,(n)}<\infty\}$, since both random variables $N_t^{\min,(n)}$ and $N_t^{\lambda,(n)}$ are integer-valued, there exists almost surely a small enough $\lambda_0>0$ such that for all $\lambda<\lambda_0$, $N_t^{\lambda,(n)}=N_t^{\min,(n)}$. Thus by the dominated convergence theorem, the upper bound~\eqref{upperboundunif} tends to $0$ as $\lambda$ goes to~$0$. Finally
\[
\underset{\lambda\,\rightarrow\,0}{\limsup} \underset{x\,\in\,[0,1]}{\sup}\bbE\left[\left(x^{N^{\min,(n)}_{t}}-x^{N_t^{\lambda,(n)}}\right)1_{\left\{N_t^{\min,(n)}\,\leq\,N_t^{\lambda,(n)}\,<\,\infty\right\}}\right]=0,
\]
and the convergence in~\eqref{unifconvsemigroup} is established. To see that the uniform convergence holds for any function $f\in C([0,1])$, one argues by the Stone--Weierstrass theorem as follows. Let $f\in C([0,1])$ and $(f_n)$ be a sequence of polynomial functions such that $\|f_n-f\|_\infty\underset{n\,\rightarrow\,\infty}{\longrightarrow} 0$. Then,
\begin{align*}
\left\|P_t^{\lambda}f-P_t^{\rmr}f\right\|_\infty &\leq \left\|P_t^{\lambda}f-P_t^{\lambda}f_n\right\|_\infty+\left\|P_t^{\lambda}f_n-P_t^{\rmr}f_n\right\|_\infty+\left\|P_t^{\rmr}f_n-P_t^{\rmr}f_n\right\|_\infty\\
&\leq 2\|f-f_n\|_\infty+\left\|P_t^{\lambda}f_n-P_t^{\rmr}f_n\right\|_\infty.
\end{align*}
By letting $\lambda$ go towards $0$, we see that
\[
\underset{\lambda\,\rightarrow\,0+}{\limsup} \left\|P_t^{\lambda}f-P_t^{\rmr}f\right\|_\infty\leq 2\|f-f_n\|_\infty
\]
and one concludes by letting $n$ go to $\infty$. We now deduce that $(P_t^{\rmr})$ is a Feller semigroup. As previously, the Stone--Weierstrass theorem asserts that it suffices to establish the semigroup property for the functions $g_n$. Let $f:=P_s^{\rmr}g_n$. For any $n\geq 0$,
\begin{multline}\label{semigroupproperty}
\left\|P_{t+s}^{\rmr}g_n-P_{t}^{\rmr}P_{s}^{\rmr}g_n\right\|_\infty \\
\begin{aligned}
&\leq \left\|P_{t+s}^{\rmr}g_n-P_{t+s}^{\lambda}g_n\right\|_\infty +\left\|P_{t}^{\lambda}P_{s}^{\lambda}g_n-P_{t}^{\lambda}P_{s}^{\rmr}g_n\right\|_\infty+\left\|P_{t}^{\lambda}P_s^{\rmr}f-P_{t}^{\rmr}P_s^{\rmr}f\right\|_\infty \\
&\leq \left\|P_{t+s}^{\rmr}g_n-P_{t+s}^{\lambda}g_n\right\|_\infty+\left\|P_{s}^{\lambda}g_n-P_{s}^{\rmr}g_n\right\|_\infty+\left\|P_{t}^{\lambda}f-P_{t}^{\rmr}f\right\|_\infty,
\end{aligned}
\end{multline}
where we have used the fact that $P_t^{\lambda} $ is a contraction. The upper bound in~\eqref{semigroupproperty} goes towards $0$ as $\lambda$ goes to $0$, and the semigroup property is established. The Feller property follows from the same argument as in the proof of Lemma~\ref{recextenslambda}. The fact that the convergence of the sequence $(X_t^{\lambda},t\geq 0,\lambda>0)$ as $\lambda$ goes to $0$, holds in the Skorokhod sense is a direct application of~\cite[Theorem~2.5 page~167]{EthierKurtz}.

\noindent\eqref{proof3} Recall the expression of $\clA^{\rms}$, the generator of the minimal solution to~\eqref{SDEselec}, given in~\eqref{generatorWFs}. Denote by $\clA^{\lambda,\rms}$ the generator of $(X_t^{\lambda},0\leq t\leq \tau)$. By Lemma~\ref{recextenslambda}, for any $g\in C_c([0,1])$, the process $(M_t^{\lambda},t\geq 0)$ defined by
\[
M_t^{\lambda}=g(X_t^{\lambda})-\int_{0}^{t}\clA^{\lambda, \rms}g(X_s^{\lambda})\ddr s
\]
is a martingale. We establish now that $\clA^{\lambda, \rms}$ converges uniformly towards $\clA^{\rms}$ on functions with a compact support $[a,b]$ contained in $(0,1)$. Recall the drift term of $\clA^{\rms}g(x)$ in~\eqref{generatorWFs}: $\mu(\bbN)x(f(x)-1)g'(x)$. Since $\clA^{\lambda, \rms}$ and $\clA^{\rms}$ have the same jump parts, for any $g\in C_c([0,1])$;
\begin{multline}\label{unifconvgen}
\left\|\clA^{\lambda, \rms}g-\clA^{\rms}g\right\|_\infty\\
\begin{aligned}
&=\underset{x\,\in\,(0,1)}{\sup}\left\lvert (\mu(\bbN)+\lambda)x\left(f^{\lambda}(x)-1\right)g'(x)-\mu(\bbN)x(f(x)-1)g'(x)\right \lvert \\
&=\underset{x\,\in\,(0,1)}{\sup}x|g'(x)|\left(\mu(\bbN)\left\lvert f^{\lambda}(x)-f(x)\right\lvert+\lambda f^{\lambda}(x)\right)\\
&\leq \mu(\bbN)\|g'\|_\infty \underset{x\,\in\,[a,b]}{\sup}\left|f(x)-f^{\lambda}(x)\right|+\lambda \|g'\|_\infty.
\end{aligned}
\end{multline}
One has
\[
\underset{x\,\in\,[a,b]}{\sup}\left|f(x)-f^{\lambda}(x)\right|=\underset{x\,\in\,[a,b]}{\sup}\sum_{k=1}^{\infty}x^{k}\mu(k)\frac{\lambda}{\mu(\bbN)+\lambda}\leq \lambda \sum_{k=1}^{\infty}\frac{\mu(k)}{\mu(\bbN)+\lambda}\leq \lambda.
\]
Therefore the right-hand side of~\eqref{unifconvgen} goes to $0$ when $\lambda$ goes to $0$ and the generators $\clA^{\lambda, \rms}$ uniformly converge to $\clA^{\rms}$. The same arguments as in the proof of Lemma~\ref{identificationXinfinity} show that
\[
\left(g\left(X^{\rmr}_{t\,\wedge\,\tau^{\rmr}}\right)-\int_{0}^{t}\clA^{\rms}g\left(X_{s\wedge \tau^{\rmr}}^{\rmr}\right)\ddr s, t\geq 0\right)
\]
is a martingale where $\tau^{\rmr}:=\inf\{t>0: X^{\rmr}_t\notin (0,1)\}$. Finally, the process $(X^{\rmr}_{t},t\geq 0)$, stopped at time $\tau^{\rmr}$ solves the martingale problem $(\MP)$. Since the latter has a unique solution, we see that $(X^{\rmr}_t,t\geq 0)$ extends the minimal process.
\end{proof}

The next lemma completes the convergence result in Lemma~\ref{recextensX} for the sequence of processes $(X^{\lambda}_t(x),t\geq 0)$ as $\lambda$ goes to $0$.
\begin{lemm}[Almost sure pointwise convergence and monotonicity]\label{almostsureinx}
For any $x\in [0,1]$ and $t\geq 0$, the limit $X^{\rmr}_t(x):=\underset{\lambda \rightarrow 0+}{\lim}\! \uparrow X_t^{\lambda}(x)$ exists almost surely. Moreover, if $x\leq y$ then $X_t^{\rmr}(x)\leq X_t^{\rmr}(y)$ a.s. In particular, the limit
\[
X_t^{\rmr}(1):=\underset{x\,\rightarrow\,1-}{\lim}\! \uparrow X_t^{\rmr}(x)\in [0,1]
\]
exists almost surely.
\end{lemm}

\begin{proof}
For any $m\geq 2$. Let $f^{\lambda}_m$ be the generating function associated to the measure $\mu^{\lambda}_m$ defined by $\mu^{\lambda}_m(k):=\mu(k)$ for $k\leq m-1$ and $\mu^{\lambda}_m(m):=\bar\mu(m)+\lambda$ for $k\geq m$. Denote by $(X_t^{\lambda,(m)},t\geq 0)$ the $\Lambda$-WF process with selection driven by $f^{\lambda}_m$. By Lemma~\ref{identificationXinfinity}, for any $x\in [0,1)$, $X_t^{\lambda}(x):=\underset{m\,\rightarrow\,\infty}{\lim} X_t^{\lambda,(m)}(x)$ and since $x^{m}-x\leq 0$,\linebreak we easily check that if $\lambda'>\lambda$,
\begin{align*}
(\mu(\bbN)+\lambda)\left(f_{m}^{\lambda}(x)-1\right)&=\sum_{k=1}^{m-1}(x^k-1)\mu(k)+(\bar{\mu}(m)+\lambda)(x^{m}-1)\\
&\geq \sum_{k=1}^{m-1}(x^k-1)\mu(k)+(\bar{\mu}(m)+\lambda')(x^{m}-1)\\
&=(\mu(\bbN)+\lambda')\left(f_{m}^{\lambda'}(x)-1\right).
\end{align*}
For any $m\geq 2$ and any $\lambda>0$, the function $f_{m}^{\lambda}$ is Lipschitz over $[0,1]$, the comparison theorem therefore applies and for any $x\in [0,1)$, $X_t^{\lambda,(m)}(x)\leq X_t^{\lambda',(m)}(x)$ a.s. By letting $m$ go to $\infty$, we get $X_t^{\lambda}(x)\leq X_t^{\lambda'}(x)$ a.s. Recall that $(X^{\lambda}_t,t\geq 0)$ can be started from $1$. By letting $x$ go to $1$, we also get $X_t^{\lambda}(1)\leq X_t^{\lambda'}(1)$ a.s. Finally, the limit $\underset{\lambda\,\rightarrow\,0+}{\lim} X_t^{\lambda}(x)=:X_t^{\rmr}(x)$ exists almost surely for all $x\in [0,1]$. The monotonicity in the initial values can be checked similarly.
\end{proof}\let\qed\relax
\end{step}
\vspace*{4pt}
%\textbf{Step 4}:
\begin{step}
We can now use both duality relationships~\eqref{dualEFC} and~\eqref{dualityEFCref} in order to classify the boundaries as in Table~\ref{correspondance} and establish Theorem~\ref{thmmomentduality2}.

\begin{lemm}\label{correspond}
The boundary $1$ is non-absorbing (respectively, inaccessible) for $(X_t^{\rmr},t\geq 0)$ if and only if the boundary $\infty$ is accessible (respectively, absorbing) for $(N_t,t\geq 0)$.
\end{lemm}
\pagebreak
\begin{proof}
Recall that by Lemma~\ref{recextensX}, $(X_t^{\rmr},t<\tau_1)$ has the same law as $(X^{\min}_t,t<\tau_1)$. As we shall use it repeatedly, we recall the duality relationships~\eqref{dualEFC} and~\eqref{dualityEFCref}: for any $n\in \bbN$, $x\in [0,1]$ and $t\geq 0$
\[
\bbE\left[x^{N^{(n)}_t}\right]\overset{\eqref{dualEFC}}{=}\bbE\left[X^{\min}_t(x)^n\right] \quad\text{and}\quad\bbE\left[x^{N^{\min,(n)}_t}\right]\overset{\eqref{dualityEFCref}}{=}\bbE\left[X^{\rmr}_t(x)^n\right].
\]
By Lemma~\ref{almostsureinx}, the limit $X^{\rmr}_t(1):=\underset{x\,\rightarrow\,1-}{\lim} X^{\rmr}_t(x)$ exists almost surely. By letting $x$ go towards $1$ in the identity~\eqref{dualityEFCref} above with $n=1$, we have $\bbE[X^{\rmr}_t(1)]=\bbP_n(\zeta_{\infty}>t)$.

For the first implication, we look at the contraposition and verify that if $\infty$ is inaccessible for $(N^{(n)}_t,t\geq 0)$, then the boundary $1$ is absorbing for $(X_t^{\rmr},t\geq 0)$. If $\infty$ is inaccessible for $(N^{(n)}_t,t\geq 0)$ then it is inaccessible for $(N^{\min,(n)}_t,t\geq 0)$ and $\bbP_1(\zeta_{\infty}>t)=\bbE(X^{\rmr}_t(1))=1$. Therefore, $\bbE(1-X^{\rmr}_t(1))=0$ and since $X^{\rmr}_t(1)\leq 1$ a.s, we get $X^{\rmr}_t(1)=1$ a.s. Thus, the boundary $1$ is absorbing for $(X_t^{\rmr},t\geq 0)$.

We then show the second implication. If $\infty$ is accessible for $(N^{(n)}_t,t\geq 0)$, then it is accessible and absorbing for $(N^{\min,(n)}_t,t\geq 0)$ and there exists $t>0$ such that $\bbP_1(\zeta_{\infty}>t)=\bbE[X^{\rmr}_t(1)]<1$. Therefore, $\bbP(X_t^{\rmr}(1)<1)>0$ and the boundary $1$ is non-absorbing for $(X_t^{\rmr},t\geq 0)$. We thus have established that $(X_t^{\rmr},t\geq 0)$ has boundary $1$ non-absorbing if and only if $\infty$ is accessible for $(N_t^{(n)},t\geq 0)$.

The second equivalence is shown along similar arguments. Letting $n$ go to $\infty$ in the identity~\eqref{dualEFC}, we get for any $x\in [0,1)$, $\bbE[x^{N^{(\infty)}_t}]=\bbP(X^{\min}_t(x)=1)$. We see that the boundary $1$ is inaccessible for the process $(X^{\min}_t,t\geq 0)$, which is equivalent to be inaccessible for $(X_t^{\rmr},t\geq 0)$, if and only if $\bbE[x^{N^{(\infty)}_t}]=0$ for any $x\in [0,1)$, which is equivalent to $N^{(\infty)}_t=\infty$ almost surely, that is to say, the boundary $\infty$ is absorbing for the process $(N_t^{(n)},t\geq 0)$.
\end{proof}\let\qed\relax
\end{step}

\begin{proof}[{Proof of Theorem~\ref{thmmomentduality2}}] The moment duality relationship~\eqref{dualityEFCref} is provided by Lemma~\ref{recextensX}. Statements~\eqref{theo2.7.i} to~\eqref{theo2.7.iv}  will be obtained by applying Lemma~\ref{correspond} and combining the necessary and sufficient conditions for boundaries to be respectively absorbing, non-absorbing and inaccessible or accessible. We provide details for statements~\eqref{theo2.7.i}, \eqref{theo2.7.ii} and~\eqref{theo2.7.iv}. Statement~\eqref{theo2.7.iii} is obtained by symmetric arguments to~\eqref{theo2.7.i}.

For statement~\eqref{theo2.7.i}, if $(N_t,t\geq 0)$ has $\infty$ as an exit boundary, then $\infty$ is absorbing and accessible. By Lemma~\ref{correspond}, $(X_t^{\rmr},t\geq 0)$ has boundary $1$ inaccessible and non-absorbing. In particular, there is $t>0$ such that $\bbE[X_t^{\rmr}(1)^{n}]=\bbP(\zeta^{(n)}_\infty>t)<1$. Hence, $\bbP(X^{\rmr}_t(1)<1)>0$. If $\tau^1$ denotes the first entrance time in $[0,1)$, then $\bbP_1(\tau^1>t)<1$. It remains to establish that almost surely $X^{\rmr}_t(1)<1$ for some $t>0$, or equivalently $\bbP_1(\tau^1=\infty)=0$. By using the Markov property at time $t$, for any $n\geq 1$,
\begin{align*}
\bbP_1\left(\tau^1>nt\right) &=\bbE_1\left(\bbmun_{\left\{\tau^1\,>\,t\right\}}\bbP_{X^{\rmr}_{t}(1)}\left(\tau^1\,>\,(n-1)t\right)\right)\\
&=\bbP_{1}\left(\tau^1>t\right)\bbP_1\left(\tau^1>(n-1)t\right).
\end{align*}
Thus, for any $n\geq 1$, $\bbP_1(\tau^1>nt)\leq \bbP_1(\tau^1>t)^n$. Since $\bbP_1(\tau^1>t)<1$, the upper bound goes to $0$ when $n$ goes to $\infty$ and we get $\bbP_1(\tau^1=\infty)=0$.

\pagebreak
For statement~\eqref{theo2.7.ii}, if $(N^{\min}_t,t\geq 0)$ has $\infty$ as a regular absorbing boundary, then the boundary $\infty$ is non-absorbing and accessible for the non-stopped process $(N_t^{(n)},\linebreak t\geq 0)$. Therefore, by Lemma~\ref{correspond}, $1$ is accessible and non-absorbing for the process $(X_t^{\rmr},t\geq 0)$, hence there is $t$ such that $\bbP(X^{\rmr}_t(1)<1)>0$. One can show as previously that almost surely $X^{\rmr}_t(1)<1$ for some $t>0$ so that $1$ is regular non-absorbing.

For statement~\eqref{theo2.7.iv}, if $(N_t,t\geq 0)$ has $\infty$ as a natural boundary, then $(N_t,t\geq 0)$ and $(N^{\min}_t,t\geq 0)$ have the same law and $\infty$ is inaccessible. The boundary $1$ of $(X_t^{\rmr},t\geq 0)$ is thus absorbing. With the boundary $\infty$ being absorbing, the boundary $1$ is also inaccessible. Hence, $1$ is a natural boundary.
\end{proof}

In the next section, we establish the correspondences given in Table~\ref{correspondancereg} and show that when boundary $1$ is regular non-absorbing, the extended $\Lambda$-WF process with selection $(X_t^{\rmr},t\geq 0)$ gets absorbed at $0$ in finite time almost-surely.


\section{Proofs of Theorems~\ref{regularforitself}, \ref{fixationat0} and Proposition~\ref{instantanprop}}\label{sec3}


\subsection{Proofs of Theorem~\ref{regularforitself} and Proposition~\ref{instantanprop}}

We show here how different boundary behaviors such as regular reflecting or regular for itself are linked by the duality relationships~\eqref{dualEFC} and~\eqref{dualityEFCref}.


\subsubsection{Proof of Theorem~\ref{regularforitself}}

\begin{proof}
Recall $\tau_1$ the first return time to $1$, and that by definition, $1$ is regular for itself if for any $t>0$, $\bbP_1(\tau_1>t)=0$. We first observe that this is equivalent to the condition $\underset{x\,\rightarrow\,1-}{\lim}\ \bbP_x(\tau_1>t)=0$. By applying the Markov property at a time $s>0$, we obtain that for any time $t>0$,
\[
\bbP_1 (\tau_1>t+s)=\bbE_1\left[\bbP_{X^{\rmr}_s(1)}(\tau_1>t)\bbmun_{\left\{\tau_1>s\right\}}\right].
\]
By the right-continuity at $0$ of the sample paths, $X^{\rmr}_s(1)\underset{s\,\rightarrow\,0+}{\longrightarrow} 1$ almost surely. Thus, $1$ is regular for itself if and only if $\underset{x\,\rightarrow\,1-}{\lim}\ \bbP_x(\tau_1>t)=0$ for all $t>0$. Note that for any $x\in [0,1)$ under $\bbP_x$,
$(X^{\rmr}_{t\wedge \tau_1}(x),t\geq 0)$ has the same law as $(X_t^{\min}(x),t\geq 0)$. By the duality relationship~\eqref{dualEFC}, for any $t>0$, $\bbP_x(\tau_1\leq t)=\bbP(X_t^{\min}(x)=1)=\bbE[x^{N_t^{(\infty)}}].$ Hence, 
\[
\underset{x\,\rightarrow\,1-}{\lim} \ \bbP_x(\tau_1\leq t)=\bbP\left(N_t^{(\infty)}<\infty\right).
\]
If the boundary $1$ of the process $(X_t^{\rmr},t\geq 0)$ is regular for itself then 
\[
\underset{x\,\rightarrow\,1-}{\lim}\ \bbP_x(\tau_1\leq t)=1\quad\text{and}\quad\bbP\left(N_t^{(\infty)}<\infty\right)=1.
\]
By Fubini's theorem, this implies that the set $\overline{\{t\geq 0: N^{(\infty)}_t=\infty\}}$ has zero Lebesgue measure almost surely, namely, $\infty$ is regular reflecting. If now $\infty$ is regular reflecting, then $\bbP(N_t^{(\infty)}<\infty)=1$ for all $t>0$, therefore $\underset{x\,\rightarrow\,1-}{\lim}\ \bbP_x(\tau_1\leq t)=1$, and as noticed before this entails that $\bbP_1(\tau_1\leq t)=1$ for all $t>0$. Hence, $\tau_1=0$, $\bbP_1$-almost surely and $1$ is regular for itself.

We now show that $1$ is regular reflecting if and only if $\infty$ is regular for itself. Set $\zeta_\infty^{(n)}:=\inf\{t>0: N_{t-}^{(n)}=\infty\}$. For any $t>0$ and $n\in \bbN$, by the duality relationship~\eqref{dualabso}, 
\[
\bbP\left(\zeta_\infty^{(n)}>t\right)=\underset{x\,\rightarrow \,1-}{\lim}\ \bbE\left[x^{N_t^{\min,(n)}}\right]=\bbE\left[X_t^{\rmr}(1)^n\right].
\]
Letting $n$ go to $\infty$ yields
\[
\underset{n\,\rightarrow\,\infty}{\lim} \bbP(\zeta_\infty^{(n)}>t)=\bbP(X_t^{\rmr}(1)=1).
\]
Provided that $1$ is regular reflecting for the process $(X_t^{\rmr},t\geq 0)$, we get $\underset{n\,\rightarrow\,\infty}{\lim} \bbP(\zeta_\infty\linebreak > t)=0$ for all $t>0$, hence, $\zeta_\infty=0$, $\bbP_{\infty}$-a.s. Therefore, $\infty$ is regular for itself for $(N_t^{(\infty)},t\geq 0)$. Similarly, if $\infty$ is regular for itself, we see that $1$ is regular reflecting.
\end{proof}

Notice that when boundary $1$ is regular reflecting, then $1$ is necessarily an instantaneous point, in the sense that $\tau^{1}:=\inf\{t>0: X_t^{\rmr}(1)<1\}=0$ a.s.


\subsubsection{Proof of Proposition~\ref{instantanprop}}

\begin{proof}
Recall the notation $\zeta_\infty^{(n)}:=\inf\{t>0: N_t^{(n)}=\infty\}$ for all $n\geq 1$ and that $\infty$ is an instantaneous exit if for all $t>0$, $\underset{n\,\rightarrow\,\infty}{\lim} \bbP(\zeta_\infty^{(n)}>t)=0$. We show how the property for boundary $1$ to be an instantaneous entrance is associated to the condition on the boundary $\infty$ of the dual process $(N_t,t\geq 0)$ to be an instantaneous exit. Recall $\tau^1$ the first entrance time in $[0,1)$ of process $(X_t^{\rmr},t\geq 0)$. The argument is similar to that in the proof of Theorem~\ref{regularforitself}. By Theorem~\ref{thmmomentduality2}, for any $t\geq 0$, $\bbE[X^{\rmr}_t(1)^{n}]=\bbP(N_t^{\min,(n)}<\infty)=\bbP(\zeta^{(n)}_\infty>t)$. Since $1$ is not accessible, for any $t>0$,
\begin{align*}
\bbP_1\left(\tau^1>t\right)&=\bbP\left(X_t^{\rmr}(1)=1\right)=\underset{n\,\rightarrow\,\infty}{\lim}\bbE\left[X^{\rmr}_t(1)^{n}\right]=\underset{n\,\rightarrow\,\infty}{\lim} \bbP\left(\zeta^{(n)}_\infty>t\right)=0
\end{align*}
which allows us to conclude the first equivalence. The second is established similarly from the first duality relationship~\eqref{dualEFC}. 
\end{proof}

\begin{rema}
The condition $\underset{n\,\rightarrow\,\infty}{\lim} \bbP(\zeta_\infty^{(n)}>t)=0$ is sometimes called $t$-regularity of the boundary $\infty$, see Kolokoltsov's book~\cite[Section~6.1, page~273]{MarkovProcessesSemigroupsandGenerators}.
\end{rema}

We study now the long term behavior of the process $(X_t^{\rmr},t\geq 0)$ when the boundary $1$ is not an exit and prove Theorem~\ref{fixationat0}

\subsubsection{Proof of Theorem~\ref{fixationat0}}

\begin{proof}
By the comparison theorem, see Theorem~\ref{comparisontheorem}, for all $x\in [0,1]$ and $t\geq 0$, $X_t^{\min}(x)\leq Y_t(x)$ a.s. where $(Y_t(x),t\geq 0)$ is a $\Lambda$-Wright--Fisher process with no selection. Under the condition~\eqref{cdipsi}, the latter reaches $0$ with positive probability, and so does the process $(X_t^{\min},t\geq 0)$.

Assume $1$ is regular non-absorbing for $(X_t^{\rmr},t\geq 0)$. Consider the successive excursions out from $1$ of the process $(X_t^{\rmr},t\geq 0)$ which crosses a given level~$x<1$. Namely, set $\tau^{(0)}_1:=0$ and $ \tau^{(n)}_x:=\inf\{t>\tau^{(n-1)}_1: X_t^{\rmr}\leq x\}$ and $\tau^{(n)}_1:=\inf\{t>\tau^{(n)}_x: X_t^{\rmr}=1\}$. Then the processes
\[
\left(X^{\rmr}_{\left(t+\tau_x^{(n)}\right)\wedge \tau_1^{(n)}},t\geq 0\right)
\]
are independent and with the same law as $(X_t^{\min},t\geq 0)$ started from $X^{\rmr}_{\tau_x^{(n)}}\leq x$ a.s. By the comparison property, each process
\[
\left(X^{\rmr}_{\left(t+\tau_x^{(n)}\right)\wedge \tau_1^{(n)}},t\geq 0\right)
\]
is stochastically below a process $(X_t^{\min}(x),t\geq 0)$ and since $\bbP_x(\tau_0<\tau_1)>0$, each excursion has a positive probability to hit $0$. Therefore, there exists almost surely an excursion among the latters which attains the boundary $0$.

Assume $1$ is an entrance for $(X_t^{\rmr},t\geq 0)$. The boundary $\infty$ is therefore an exit for the process $(N^{\min}_t,t\geq 0)$ and by the duality relationship~\eqref{dualabso}, we get for all $x\in [0,1]$,
\[
\underset{t\,\rightarrow\,\infty}{\lim}\bbE\left[X_t^{\rmr}(x)\right]=\underset{t\,\rightarrow\,\infty}{\lim} \bbE\left[x^{N_t^{\min}}\right]=0.
\]
Hence $\underset{t\,\rightarrow\,\infty}{\liminf}\ X_t^{\rmr}(x)=0$ a.s. Set $\tau_{1/n}:=\inf\{t>0: X_t^{\rmr}\leq 1/n\}$. For all $n\geq 2$, $\tau_{1/n}<\infty$ a.s. Since $1$ is an entrance boundary, for any $x\in [0,1)$, $(X_t^{\rmr}(x),t\geq 0)$ has the same law as $(X_t^{\min}(x),t\geq 0)$. By the strong Markov property at time $\tau_{1/n}$, $(X_{t+\tau_{1/n}}^{\rmr}(x),t\geq 0)$ has the same law as $(X_t^{\min}(X^{\rmr}_{\tau_{1/n}}(x)),t\geq 0)$, where $(X_t^{\min},t\geq 0)$ is independent from $X^{\rmr}_{\tau_{1/n}}(x)$. Since $X^{\rmr}_{\tau_{1/n}}(x)\leq 1/n$ a.s, by the comparison theorem, $(X_{t+\tau_{1/n}}^{\rmr}(x),t\geq 0)$ is stochastically smaller than $(Y_t(1/n),t\geq 0)$ where $(Y_t(1/n),\linebreak t\geq 0)$ is a $\Lambda$-Wright--Fisher process with no selection.


Set $E_n:=\{X^{\min}_{t+\tau_{1/n}}>0, \forall t\geq 0\}$ for any $n\geq 2$. One has
\[
\bbP(E_n)\leq \bbP(Y_t(1/n)>0 \text{ for all } t\geq 0)=\bbP_{1/n}\left(\tau^Y_0>\tau^Y_1\right),
\]
where $\tau^Y_i:=\inf\{t\geq 0; Y_t=i\}$ for $i=0,1$. By the duality relationship for the pure $\Lambda$-coalescent: $\bbP_x(\tau^Y_1<\tau^Y_0)=\bbE_\infty[x^{N^{Y}_t}]\underset{x\,\rightarrow\,0}{\longrightarrow} 0$ where we have denoted by $(N_t^{Y},t\geq 0)$ the moment dual of $(Y_t,t\geq 0)$. Thus $\bbP(E_n)\underset{n\,\rightarrow\,\infty}{\longrightarrow} 0$. Since $0$ is an absorbing boundary, $E_{n+1}\subset E_n$ for all $n\geq 2$. Hence $\bbP\left(\cap_{n=2}^{\infty}E_n\right)=\underset{n\,\rightarrow\,\infty}{\lim} \bbP(E_n)=0.$ This allows us to conclude since $\cup_{n=2}^{\infty} E_n^{c}$ has probability $1$ and
\[
\cup_{n=2}^{\infty} E_n^{c}\subset \left\{\exists t_0\geq 0; X_t^{\rmr}(x)=0, \ \forall t\geq t_0\right\}.\qedhere
\]
\end{proof}

Until now we have only shown theoretical results on possible extensions of the minimal process and their duality relationships with the process $(N_t,t\geq 0)$ and the stopped process $(N^{\min}_{t},t\geq 0)$. One may wonder whether there exist mechanisms of resampling $\Lambda$ and selection $f$ that result in the regular boundary, see the last line in Table~\ref{correspondance}. It is not clear whether the easiest route to study a given process is to look at its dual or not. The aim of the next section is to translate all results known about block counting processes for simple EFCs recalled in Section~\ref{backgroundEFC} into results for the $\Lambda$-Wright--Fisher process with selection.


\section{Proofs of Theorems~\ref{theorementrancef}, \ref{suffcondexitf}, \ref{regularforitselfstable} and Corollary~\ref{inftyregularreflecting}}\label{application}


We transfer the results known for the block counting process $(N_t,t\geq 0)$, see Section~\ref{backgroundEFC}, to the $\Lambda$-Wright--Fisher process with frequency-dependent selection by applying our two duality relationships~\eqref{dualEFC} and~\eqref{dualityEFCref}. Recall the correspondences stated in Table~\ref{correspondance} and Table~\ref{correspondancereg}. In the sequel, we work with the extension $(X_t^{\rmr},t\geq 0)$ of the minimal solution to~\eqref{SDEselec}, $(X_t^{\min},t< \tau)$, which is constructed in Lemma~\ref{recextensX}.

\subsection{Proof of Theorem~\ref{suffcondexitf}} 

\begin{proof}
Recall Lemma~\ref{nonexplosioncriterion}. If $\sum_{n=2}^{\infty}\frac{n}{\Phi(n)}\bar{\mu}(n)<\infty$, then the block counting process $(N_t,t\geq 0)$ does not explode, i.e. $\infty$ is inaccessible. In this case, by Lemma~\ref{correspond}, the extended process $(X_t^{\rmr},t\geq 0)$ has its boundary $1$ absorbing. If moreover, $\sum_{n=2}^{\infty}\frac{1}{\Phi(n)}\linebreak<\infty$ (which is equivalent to $\int^{\infty}\frac{\ddr x}{\Phi(x)}<\infty$), then the process $(N_t,t\geq 0)$ has boundary $\infty$ as an entrance. Recall that~\cite[Lemma~2.5]{cdiEFC} guarantees that $\infty$ is instantaneous. Theorem~\ref{thmmomentduality2} and Proposition~\ref{instantanprop} ensure then that $1$ is an instantaneous exit for $(X_t^{\rmr},t\geq 0)$. Last, if $\sum_{n=2}^{\infty}\frac{1}{\Phi(n)}=\infty$, then $\infty$ is a natural boundary and by Theorem~\ref{thmmomentduality2}, $1$ is also natural. Therefore it only remains to show that the assumptions of Theorem~\ref{suffcondexitf} entail the convergence of the series $\sum_{n=2}^{\infty}\frac{n}{\Phi(n)}\bar{\mu}(n)$.

Assume first that $f$ is Lipschitz on $[0,1]$, namely $f'(1-)<\infty$. The splitting measure $\mu$ admits then a first moment. Recall that the sequence $(\Phi(n)/n, n\geq 2)$ is nondecreasing, see e.g.~\cite[Lemma~2.1-(iv)]{limic2015}. Hence the sequence $(\frac{n}{\Phi(n)},n\geq 2)$ is bounded and $\sum_{n=2}^{\infty}\frac{n}{\Phi(n)}\bar{\mu}(n)\leq \frac{2}{\Phi(2)}\sum_{n=2}^{\infty}\bar{\mu}(n)<\infty$. This establishes the first case.

Assume now that the selection function $f$ satisfies some properties of regular variation. In this setting, we can use Tauberian theorems to relate the asymptotics of the selection function $f$ at boundary $1$, with asymptotics of the splitting measure $\mu$ at $\infty$. We gather here these results.

Set $u(x):=\mu(\bbN)(1-f(x))$ for all $x\in [0,1]$ and recall $\ell(n)=\sum_{k=1}^{n}\bar{\mu}(k)$. For all $\lambda \geq 0$, set $\kappa(\lambda):=\int_{0}^{\infty}(1-e^{-\lambda x})\mu(\ddr x)$ with $\mu(\ddr x)=\sum_{k=1}^{\infty}\mu(k)\delta_k$. Thus,
\begin{equation}\label{equiv}
u(e^{-\lambda})=\mu(\bbN)\left(1-f(e^{-\lambda})\right)=\sum_{k=1}^{\infty}\left(1-e^{-\lambda k}\right)\mu(k)=\kappa(\lambda).
\end{equation}

Let $s$ be a slowly varying function at $\infty$. By the Tauberian theorem, see e.g.~\cite[Chapter~0.7 page~10]{Bertoin96}, the following are equivalent:
\begin{enumerate}\romanenumi
\item\label{sect7.1.i} $\mu(n)\underset{n\,\rightarrow\,\infty}{\sim} \frac{b }{n^{\alpha+1}}s(n)$ for some $\alpha\in (0,1)$,
\item\label{sect7.1.ii} $\kappa(\lambda)\underset{\lambda\,\rightarrow\,0+}{\sim} \lambda^{\alpha}\frac{b\Gamma(2-\alpha)}{\alpha(1-\alpha)}s(1/\lambda)$,
\item\label{sect7.1.iii} $\mu(\bbN)(1-f(x))=\kappa(\log 1/x)\underset{x\,\rightarrow\,1-}{\sim} \frac{b\Gamma(1-\alpha)}{\alpha}(1-x)^{\alpha}s(\frac{1}{1-x}).$
\end{enumerate}
Similarly, if $s$ is slowly varying at $\infty$, then we have the equivalence:
\begin{enumerate}
\item $\mu(\bbN)(1-f(x))\underset{x\,\rightarrow\,1-}{\sim}\kappa(\log 1/x)\underset{x\,\rightarrow\,1-}{\sim} s(1/(1-x))$,
\item $\bar{\mu}(n)\underset{n\,\rightarrow\,\infty}{\sim} s(n)$.
\end{enumerate}
Recall the Tauberian equivalence~\eqref{sect7.1.iii} $\Longleftrightarrow$\eqref{sect7.1.i}. When $x\mapsto 1-f(x)$ is regularly varying at $1$ with index $\alpha\in [0,1)$, $1-f(x)\underset{x\,\rightarrow\,1-}{\sim} c(1-x)^{\alpha}s(1/(1-x))$ for some slowly varying function $s$ at $\infty$ and $\bar{\mu}(n)\underset{n\,\rightarrow\,\infty}{\sim} \frac{c'}{n^{\alpha}}s(n)$ for some constant $c'>0$. Hence, $\bar{\mu}(\frac{1}{1-x})\underset{x\,\rightarrow\,1-}{\sim} c'(1-x)^{\alpha}s(1/(1-x)) \underset{x\,\rightarrow\,1-}{\sim} c'' (1-f(x))$ where $c''$ is a positive constant. Simple integral comparisons and a change of variable give the equivalences:
\begin{align*}
\sum_{n=2}^{\infty}\frac{n}{\Phi(n)}\bar{\mu}(n)<\infty &\Longleftrightarrow \int^{1-}\frac{1/(1-x)}{\Phi\big(1/(1-x)\big)} (1-x)^{\alpha}s\big(1/(1-x)\big)\frac{\ddr x}{(1-x)^2}<\infty \\
&\Longleftrightarrow \int^{1-}\frac{1-f(x)}{(1-x)^3\Phi\big(1/(1-x)\big)} \ddr x<\infty.
\end{align*}
The integral condition above matches with our condition~\eqref{integralform2}.
\end{proof}


\subsection{Proof of Theorem~\ref{theorementrancef}}


\begin{proof}
Recall the assumptions $\Lambda(\{0\})=\Lambda(\{1\})=0$. According to Lemma~\ref{recextenslambda}, when the function $f$ is defective, namely $f(1)<1$, the $\Lambda$-WF process with selection $(X_t,t\geq 0)$, minimal solution to~\eqref{SDEselec}, has boundary $1$ entrance. We are interested here on the non-defective selection functions for which $f(1)=1$.


If the boundary $\infty$ of the process $(N_t,t\geq 0)$ is an exit, then Theorem~\ref{thmmomentduality2} implies that the process $(X_t^{\rmr},t\geq 0)$ has boundary $1$ entrance. The fact that $1$ is an entrance boundary will therefore be a simple consequence of Lemma~\ref{suffcondexitN}. We first establish that condition $\clH$ for the function $x\mapsto 1-f(x)$ entails the condition $\bbH$ for the function $\ell :n\mapsto \sum_{k=1}^{n}\bar{\mu}(k)$.

Recall $u(x):=\mu(\bbN)(1-f(x))$ for any $x\in (0,1)$ and the identity~\eqref{equiv}. By~\cite[Proposition~1, Chapter~III]{Bertoin96}, there is a universal constant $c>1$, such that
\[
\frac{1}{c}\int_{1}^{1/\lambda}\bar{\mu}(x)\ddr x \leq \frac{\kappa(\lambda)}{\lambda}\leq c \int_{1}^{1/\lambda}\bar{\mu}(x)\ddr x,
\]
where $\bar{\mu}(x)=\bar{\mu}(k)$ for any $x\in [k,k+1[$. One can check $\int_{1}^{1/\lambda}\bar{\mu}(x)\ddr x \underset{\lambda\,\rightarrow\,0}{\sim} \ell(\lfloor 1/\lambda \rfloor).$ By change of variable $\lambda=\log 1/x$, the equivalence just mentioned and the identity~\eqref{equiv}, we see that for $x$ close enough to $1$, then
\begin{equation}\label{comparisonlu}
\underbar c \ell\left(\left \lfloor 1/\log 1/x \right \rfloor \right)\leq \frac{u(x)}{\log 1/x}\leq \bar c \ell\left(\left \lfloor 1/\log 1/x \right \rfloor \right)
\end{equation}
for some constants $\bar c>c>\underbar c $. We now show that condition $\clH$ entails condition $\bbH$.

By condition $\clH$, if $x$ is close enough to $1$, then $u(x)\geq \mu(\bbN)L(x)$ for some function $L$ such that the map $h$ satisfying $\mu(\bbN)L(x)=(1-x)\log\left(1/(1-x)\right) h(x)$, is non-decreasing. Moreover, since $1-x \underset{x\rightarrow 1-}{\sim} \log 1/x$, if $x$ is close enough to $1$, then
\[
u(x)\geq \mu(\bbN)L(x)=(1-x)\log\left(1/(1-x)\right) h(x) \geq C(\log 1/x) (\log 1/\log 1/x) h(x)
\]
for some constant $C>0$. By applying the upper bound in~\eqref{comparisonlu} in the inequality above, we see that
\[
\ell\left(\lfloor 1/\log 1/x \rfloor \right)\geq C\log(1/\log 1/x) h(x),
\]
for some other constant $C>0$. Thus, for large enough $n$,
$\ell(n)\geq C(\log n)g(\log n)$ with $g$ the map such that $g(\log 1/\log 1/x):=h(x)$. By assumption, the map $h$ is non-decreasing in a neighbourhood of $1$, the map $g$ is therefore eventually non-decreasing. One also easily checks that $\int^{1}\frac{1}{L(x)}\ddr x<\infty$ entails $\int^{\infty}\frac{1}{xg(x)}\ddr x<\infty$. Finally, Condition $\bbH$ holds.

Under the condition~\eqref{entrancecondu}, we see from the upper bound in~\eqref{comparisonlu}, that $\underset{n\,\rightarrow\,\infty}{\lim}\frac{\Phi(n)}{n\ell(n)}=0$, hence Lemma~\ref{suffcondexitN} applies and $1$ is an entrance. It remains to show that the entrance at boundary $1$ is instantaneous. It has been established in~\cite[Proposition~5.2]{explosion}, that under condition $\bbH$, when $\underset{n\,\rightarrow\,\infty}{\lim} \frac{\Phi(n)}{n\ell(n)}=0$, the boundary $\infty$ of process $(N_t^{(n)},t\geq 0)$ is an instantaneous exit. We can therefore apply Proposition~\ref{instantanprop} which ensures that boundary $1$ of the dual process $(X_t^{\rmr},t\geq 0)$ is an instantaneous entrance.
\end{proof}


\begin{rema}\label{tauberianremark}
Note that the inequalities~\eqref{comparisonlu} entail that
$\int^{1-}\frac{\ddr x}{1-f(x)}<\infty$ if and only if $\sum_{n\,\geq\,1}\frac{1}{n\ell(n)}<\infty$. We recover analytically here the equivalence between Dynkin's condition and Doney's condition for explosion of a pure discrete branching process in continuous time whose offspring law is $\mu$ and generating function of $\mu$ is $f$, see Doney~\cite{Doney1984} and e.g.~\cite[Section~2.4]{explosion}.
\end{rema}

\subsection{Proofs of Theorem~\ref{regularcase}, Theorem~\ref{regularforitselfstable} and Corollary~\ref{inftyregularreflecting}}

\subsubsection{Proof of Theorem~\ref{regularcase}}
\begin{proof}
Tauberian theorems ensure that the conditions over $\Lambda$ and $f$ are equivalent to
\[
\Phi(n)\underset{n\,\rightarrow\,\infty}{\sim} d n^{\beta+1}
\text{ and }
\mu(n)\underset{n\,\rightarrow\,\infty}{\sim} \frac{b}{n^{1+\alpha}}\text{ with }d:=\frac{\Gamma(1-\beta)}{\beta(1+\beta)}\rho \text{ and } b:=\frac{\alpha}{\Gamma(1-\alpha)}\sigma,
\]
see e.g.~\cite[Section~2.2, p.~12]{cdiEFC} for the first equivalence. Cases~\eqref{theo2.17.i} and~\eqref{theo2.17.ii} are obtained by applying respectively Theorem~\ref{theorementrancef} and Theorem~\ref{suffcondexitf} previously established. Lemma~\ref{stablefragtheorem} classifies the boundary $\infty$ when $\alpha=1-\beta$ according to the ratio $b/d$. The cases: entrance, exit and regular in iii) are obtained by noticing that
\[
\sigma/\rho>\frac{\pi}{(2-\alpha)\sin(\pi \alpha)}\text{ is equivalent to }\frac{b}{d}>\alpha(1-\alpha)\text{ and }\sigma/\rho>\frac{1}{(1-\alpha)(2-\alpha)}
\]
is equivalent to $b/d>\frac{\alpha \sin(\pi \alpha)}{\pi}$. In particular, when $\frac{1}{(1-\alpha)(2-\alpha)}<\sigma/\rho<\frac{\pi}{(2-\alpha)\sin(\pi \alpha)}$, the process $(N_t,t\geq 0)$ has $\infty$ as a regular non-absorbing boundary. Therefore, the process $(N^{\min}_t,t\geq 0):=(N_{t\,\wedge\,\zeta_{\infty}},t\geq 0)$ has $\infty$ as a regular absorbing boundary, and Theorem~\ref{thmmomentduality2} ensures that the process $(X_t^{\rmr},t\geq 0)$ has boundary $1$ regular \emph{non-absorbing}. We now need to check that the boundary $1$ is instantaneous when it is an exit or an entrance and is reflecting when it is regular. In the case of $1$ being an exit for $(X_t^{\rmr},t\geq 0)$, the boundary $\infty$ is an entrance for $(N_t,t\geq 0)$. According to~\cite[Lemma~2.5]{cdiEFC}, $\infty$ is an instantaneous entrance boundary and Proposition~\ref{instantanprop} entails then that $1$ is an instantaneous exit. To deal with the case of $1$ regular non-absorbing, one applies~\cite[Proposition~3.7]{explosion}, which ensures that the dual process $(N_t,t\geq 0)$ has boundary $\infty$ regular for itself. The fact that $1$ is reflecting is a consequence of Theorem~\ref{regularforitself}. Only remains to show that $1$ is an instantaneous entrance when $\sigma/\rho>\frac{\pi}{(2-\alpha)\sin(\pi \alpha)}$, i.e. $\frac{b}{d}>\alpha(1-\alpha)$. The argument given in the proof of~\cite[Proposition~3.7]{explosion} actually covers this case since $\frac{b}{d}>\alpha(1-\alpha)$ entails $\frac{b}{d}>\alpha \sin(\pi \alpha)/\pi$. This ensures that $\infty$ is an instantaneous exit. By applying again Proposition~\ref{instantanprop}, one gets that $1$ is an instantaneous entrance.
\end{proof}

\subsubsection{Proof of Theorem~\ref{regularforitselfstable}}

\begin{proof}
We now specify the behavior of the process $(X_t^{\rmr},t\geq 0)$ at its boundary $1$ when it is regular reflecting by showing that the boundary is regular for itself, namely $\tau_1=0$, $\bbP_1$-a.s. Similarly as in the proof of Theorem~\ref{regularforitself}, the boundary $1$ is regular for itself if and only if
\begin{equation}\label{pxtau1}
\bbP_x(\tau_1>t)\underset{x\,\rightarrow\,1-}{\longrightarrow} 0 \text{ for any } t>0.
\end{equation}
We now establish~\eqref{pxtau1}. Recall $\clA$ in~\eqref{generator} and $\clA^{\rms}$ in~\eqref{generatorWFs}. Let $\epsilon\in (0,1)$. By~\cite[Proposition~6.3.2, p.~281]{MarkovProcessesSemigroupsandGenerators}, the existence a positive function $g$ on $[0,1]$ such that $g\in C^{2}([0,1])$, $g(1)=0$ and there is $c>0$, such that $\clA^{\rms}g(x)\leq -c$ for any $x\in (1-\epsilon,1)$, entails~\eqref{pxtau1}. We now look for such a function $g$. Observe that
\begin{align*}
\clA g(x) &=  \ x\int_{0}^{1}\big(g(x+z(1-x))-g(x)-z(1-x)g'(x)\big) z^{-2}\Lambda(\ddr z)\\
&+(1-x)\int_{0}^{1}\big(g(x(1-z))-g(x)+zxg'(x)\big) z^{-2}\Lambda(\ddr z)\\
&:=\clA^{+}g(x)+\clA^{-}g(x).
\end{align*}
Recall $\clA^{\rms}g(x)=\clA g(x)+\mu(\bbN)x(f(x)-1)g'(x)$ for al $x\in [0,1]$, see~\eqref{generatorWFs}. Note that for any function $g\in C^{2}((0,1))$, one has for any $y, u$,
\[
g(y+u)-g(y)-ug'(y)=u^2\int_0^1 g''(y+vu)(1-v)\ddr v.
\]
Let $g(x)= (1-x)^\delta$ for $0<x<1$ and $0<\delta<1$. Plainly,
\[
g(1)=0,\quad g'(x)= -\delta(1-x)^{\delta-1}<0, \quad g''(x)=\delta(\delta-1)(1-x)^{\delta-2}<0.
\]
Since by assumption $\mu(n)\underset{n\,\rightarrow \,\infty}{\sim} \frac{b}{n^{1+\alpha}}$, then by Item~\eqref{sect7.1.ii} below Equation~\eqref{equiv}, we see that and
\[
1-f(x) \underset{x\,\rightarrow\,1-}{\sim} C\log(1/x)^{\alpha}.
\]
Hence, the generating function $f$ satisfies $\limsup_{x\,\rightarrow\,1-}\frac{x-f(x)}{(1-x)^\alpha}<\infty$. Choosing $\delta=1-\alpha$, we have $x(f(x)-1)g'(x) =-\delta x(f(x)-1)(1-x)^{\delta-1}<c$ for some $c>0$ and all $x$ close to $1$. In addition, for all $x\in [0,1]$,
\begin{equation}\label{A+}
\begin{split}
\clA^{+}g(x)&=x\int_0^1 (g(x+z(1-x))-g(x)-z(1-x) g'(x))z^{-2}\Lambda(\ddr z) \\
&=\delta(\delta-1)x(1-x)^2\int_0^1 \Lambda(\ddr z)\int_0^1(1-x-vz(1-x))^{\delta-2}(1-v)\ddr v\\
&=\delta(\delta-1)x(1-x)^\delta\int_0^1\Lambda(\ddr z)\int_0^1(1-vz)^{\delta-2}(1-v)\ddr v<0,
\end{split}
\end{equation}
and for all $x$ close to $1$,
\begin{equation}\label{A-}
\begin{split}
\clA^{-}g(x)&=(1-x)\int_0^1(g(x(1-z))-g(x)+zxg'(x))z^{-2}\Lambda(\ddr z) \\
&=\delta(\delta-1)(1-x)x^2\int_0^1 \Lambda(\ddr z)\int_0^1 (1-x+vxz)^{\delta-2}(1-v)\ddr v\\
&\leq c_1(\delta-1)(1-x)\int_0^1 (1-x+z)^{\delta-2} \Lambda(\ddr z)\int_0^1(1-v)\ddr v \\
&\leq c_2(\delta-1)(1-x)\int_0^{1-x} (1-x+z)^{\delta-2} \Lambda(\ddr z) \\
&\leq -c_3(1-x)(1-x)^{\delta-2}\int_0^{1-x}\Lambda(\ddr z) \leq -c_3(1-x)^{\delta-1}\int_0^{1-x}\Lambda(\ddr z),
\end{split}
\end{equation}
where $c_i, i\in \{1,2,3\}$ are positive constants. By assumption $\Lambda(\ddr z)=h(z)\ddr z$ with $h(z)\underset{z\,\rightarrow\,0}{\sim} \rho z^{\alpha-1}$. Thus, $\int_0^{1-x}\Lambda(\ddr z)\underset{x\,\rightarrow\,1-}{\sim} C(1-x)^{\alpha}$ for some constant $C>0$. Combining~\eqref{A+} and~\eqref{A-}, and recalling that $\delta=1-\alpha$, we see that
\[
\underset{x\,\rightarrow\,1-}{\limsup}\ A^sg(x)\leq -c
\] 
for some positive constant $c$.
\end{proof}


\subsubsection{Proof of Corollary~\ref{inftyregularreflecting}}
\begin{proof}
The proof is straightforward. By Theorem~\ref{regularforitselfstable}, the boundary $1$ of $(X_t^{\rmr},t\geq 0)$ is regular for itself. Then statement~\eqref{coro2.21.i} follows by applying Theorem~\ref{regularforitself}. Statement~\eqref{coro2.21.ii} is a consequence of Theorem~\ref{rec}. 
\end{proof}


Other explicit cases have been found in~\cite{explosion} and have their counterparts for the dual $\Lambda$-Wright--Fisher processes with frequency-dependent selection. We refer for instance to~\cite[Theorem~3.11]{explosion} for the case $\Phi(n)\underset{n\,\rightarrow\,\infty}{\sim} d n(\log n)^{\beta}$ and $\mu(n)\underset{n\,\rightarrow\,\infty}{\sim} b \frac{(\log n)^{\alpha}}{n^2}$. We finally mention a question that has not been addressed in the present article. The equivalence stated in Theorem~\ref{regularforitself} entails that if one of the processes has its boundary \emph{irregular for itself} then the other process has its boundary \emph{sticky}, in the sense that the process stays at the boundary a set of times of positive Lebesgue measure. More precisely, for the process $(X_t^{\rmr}(x),t\geq 0)$, the duality relation~\eqref{dualabso} implies that for any time 
\[
t\geq 0,\,\underset{n\,\rightarrow\,\infty}{\lim}\bbE\left[x^{N_t^{\min,(n)}}\right]=\bbP\left(X_t^{\rmr}(x)=1\right),
\]
which might be positive if the process $(N^{(\infty)}_t,t\geq 0)$ does not return to $\infty$ instantaneously i.e. $\underset{n\,\rightarrow\,\infty}{\lim}\bbP_n(\zeta_\infty> t)>0$ for some $t>0$. Whether there are some resampling measure $\Lambda$ and selection function $f$ for which boundaries $1$ or $\infty$ are indeed sticky or irregular for themselves remains an unsolved question. 

\subsection*{Acknowledgements}
The authors are grateful to two anonymous referees for very detailed suggestions and comments.
%\providecommand{\MR}{\relax\ifhmode\unskip\space\fi MR }
%\providecommand{\MRhref}[2]{\href{http://www.ams.org/mathscinet-getitem?mr=#1}{#2} }

\bibliography{AHL_Foucart}
\end{document}
